\documentclass[../main.tex]{subfiles}
\begin{document}

\chapter{ML Accelerations On GPUs}
\lessoninfo{
  video={M12, lessons 1--3 (GPU and ML; Floating Point Formats; Tensor Core)},
  sources={NVIDIA A100 \cite{nvidia2020ampere} and H100 \cite{nvidia2022hopper} whitepapers (required, as the tensor-core resource); NVIDIA GEMM guide \cite{nvidia2026gemmguide} (optional); PTX ISA \cite{nvidia2026ptx}; IEEE 754-2019 \cite{ieee754-2019}; Cambier et al.\ \cite{cambier2020s2fp8} (S2FP8); Micikevicius et al.\ \cite{micikevicius2022fp8} (FP8); Kung \cite{kung1982systolic}; Jouppi et al.\ \cite{jouppi2017tpu}; Dao et al.\ \cite{dao2022flashattention,dao2023flashattention2} (FlashAttention, background)},
  textbook={K12 \cite{kim2012gpgpu} §1.3.9; PMPP \cite{hwu2022pmpp} ch.~14, ch.~16, App.~A; CPG \cite{nvidia2024cudaguide} §7.24 Warp Matrix Functions, §7.26 Asynchronous Barrier, §7.27 Asynchronous Data Copies; BPG \cite{nvidia2026cudabestpractices} §7.3 Numerical Accuracy and Precision},
  keywords={machine learning, GEMM, arithmetic intensity, ops:byte ratio, tensor core, matrix multiply-accumulate, systolic array, Transformer Engine, floating-point format, IEEE 754, subnormal number, bfloat16, TF32, FP8, quantization, mixed precision, structured sparsity, SpMV, asynchronous copy},
}

Machine learning is the workload that has shaped GPUs most in the last
decade: the tensor cores, FP16, FP8 and the Transformer Engine of
Lesson~1 were all added for it.  This last lesson asks why GPUs suit
machine learning and what dedicated hardware they add for it.  It
first identifies the operations that dominate neural networks, above
all matrix multiplication, and measures them with the arithmetic
intensity and Roofline reasoning of Lesson~8.  It then follows the
number formats from IEEE single precision down to 8 bits, because the
format decides how many bytes each operation moves.  The last part
turns to the design of the matrix units themselves: which operation to
accelerate, how to build the unit, how to feed it from memory, and how
to exploit the zeros of sparse matrices.  The compiler view of the
previous two lessons stays in the background: the new units appear to
software as new instructions of the ISA.

\section{GPU and ML}
\label{sec:l12-gpu-ml}

The field of \term{machine learning (ML)}[機械学習] builds models whose
parameters are fitted to data, and its dominant model today is the
\term{deep neural network (DNN)}[深層ニューラルネットワーク], a stack of
layers that each transform a large array of numbers.\source{M12, GPU and
ML; PMPP ch.~16; \cite{nvidia2026gemmguide}.}  This section explains why
GPUs suit this workload, which operations it consists of, how
matrix multiplication is measured, and the two hardware routes to
accelerating it: SIMD lanes and dedicated matrix units.

\subsection{Why GPUs suit machine learning}

Four properties of GPUs match the needs of ML training and inference.
\begin{description}
  \item[Floating-point throughput.]  ML performs a very large number of
    floating-point operations with high data parallelism, and most of a
    GPU's transistors serve floating-point throughput (Lesson~1).
  \item[Memory bandwidth.]  The GPU and its memory are connected
    directly on the package, so stacked high bandwidth memory (HBM) is
    easy to adopt; a CPU socket must keep its standard memory
    interfaces.\margin{HBM bandwidth per GPU:
    \SI{900}{\giga\byte\per\second} (V100),
    \SI{1555}{\giga\byte\per\second} (A100 40\,GB),
    \SI{3000}{\giga\byte\per\second} (H100 SXM5, preliminary) (AWP
    Table~4 \cite{nvidia2020ampere}; HWP Table~3
    \cite{nvidia2022hopper}).}
  \item[Flexible formats.]  GPUs add new number formats quickly: tensor
    formats, bfloat16, 8-bit floating point, INT8 and INT4
    (\cref{sec:l12-formats}).
  \item[Relaxed semantics.]  ML is statistical computing: a result
    that differs in the last bits, or an operation reordered between
    threads, does not change what the model learns (floating-point
    addition is not associative, so a reordering changes only the last
    bits: BPG §7.3.2 Floating Point Math Is Not Associative).  This fits the GPU's
    weak memory consistency and its lack of precise exceptions
    (Lesson~1), both of which would cost hardware to
    provide.\margin{Lesson~1 contrasts the precise exceptions of a CPU
    with the GPU, which does not guarantee them.}
\end{description}

\subsection{Operations of a neural network}

The layers of a DNN are built from three classes of operations, all
highly data parallel.

\begin{definition}[Classes of DNN operations]
  The operations of a DNN fall into three classes.
  \begin{itemize}
    \item An \term{element-wise operation}[要素ごとの演算] applies the
      same function independently to every element of an array.  The
      typical example is an \term{activation function}[活性化関数], the
      non-linear function applied to every output of a layer, such as
      $\mathrm{ReLU}(x) = \max(0, x)$.
    \item A reduction (Lesson~1) combines many elements into one; the
      DNN example is \term{pooling}[プーリング], which replaces each
      window of an array by one value, such as its maximum (max
      pooling).
    \item A \term{dot product}[内積] $\sum_k a_k b_k$ multiplies two
      vectors element by element and sums the products.  Dot products
      make up matrix multiplications and convolutions.
  \end{itemize}
\end{definition}

The dot-product class dominates the arithmetic.  A
\term{fully connected layer}[全結合層] connects every input to every
output with its own weight, so a batch of inputs times the weight
matrix is one \term{matrix multiplication}[行列積].  A
\term{convolution layer}[畳み込み層] slides a small filter over its
input and is likewise computed as a matrix multiplication (PMPP
ch.~16).\margin{Element-wise operations and reductions read and write
every element once and are therefore memory-bound in the sense of the
Roofline model (Lesson~8); only the dot-product class reuses its
operands.}

\subsection{GEMM and its arithmetic intensity}

Libraries express the matrix multiplications of ML in one standard
form.

\begin{definition}[GEMM]
  The operation \term{GEMM (general matrix multiplication)}[GEMM（一般行列積）]
  computes
  \begin{equation}
    C \leftarrow \alpha A B + \beta C ,
    \label{eq:l12-gemm}
  \end{equation}
  where $A$ is an $M \times K$ matrix, $B$ is $K \times N$, $C$ is
  $M \times N$, and $\alpha$, $\beta$ are scalars.  With $\alpha = 1$ and
  $\beta = 0$ it is the plain product $C = AB$.
\end{definition}

Each of the $MN$ elements of $C$ is a dot product of length $K$, so a
GEMM performs $MNK$ fused multiply-adds (Lesson~8), that is $2MNK$
floating-point operations.\margin{The GEMM guide counts \enquote{a total
of \texttt{2 * M * N * K} FLOPS} \cite[§2]{nvidia2026gemmguide}.}  With
FP16 inputs (two bytes per element) it reads $A$ and $B$ and writes $C$,
$2(MK + NK + MN)$ bytes if every matrix crosses the memory interface
exactly once.  The ratio is its arithmetic intensity (Lesson~8):
\begin{equation}
  I_{\text{GEMM}}
  = \frac{2MNK}{2(MK + NK + MN)}
  = \frac{MNK}{MK + NK + MN}\ \si{\flop\per\byte} .
  \label{eq:l12-ai}
\end{equation}
\caution{The intensity of \cref{eq:l12-ai} is a first-order estimate:
it assumes each matrix is read or written exactly once from memory and
ignores caches and on-chip reuse.  A real kernel moves more bytes
(tiles of $A$ and $B$ are read several times) or fewer (operands that
stay in L2).}%
The machine has an intensity of its own, the peak arithmetic rate
divided by the memory bandwidth, and the comparison of the two decides
which resource limits the kernel.

\begin{definition}[Ops:byte ratio; math- and memory-limited]
  The \term{ops:byte ratio}[演算／バイト比] of a processor is its peak
  arithmetic rate divided by its peak memory bandwidth, in FLOP per
  byte; it is the ridge point of the Roofline model of Lesson~8.  The
  side of the ratio on which a kernel's intensity lies makes it
  \term{math-limited / memory-limited}[演算律速／メモリ律速]: above the
  ratio the arithmetic units are the bottleneck (math-limited), below it
  the memory bandwidth (memory-limited).
\end{definition}

\begin{example}
  \label{ex:l12-ai}
  A square GEMM with $M = N = K = 1024$ in FP16 has
  $I = 1024^3 / (3 \cdot 1024^2) = 1024/3 \approx
  \SI{341}{\flop\per\byte}$, well above the
  \SI{138.9}{\flop\per\byte} of a V100,\margin{V100 with FP16 tensor-core
  inputs and FP32 accumulation: \enquote{the \texttt{FLOPS:B} ratio is
  138.9 if data is loaded from the GPU's memory}
  \cite[§2]{nvidia2026gemmguide}.} so it is math-limited.  A skinny GEMM
  with $M = K = 4096$ and $N = 16$, a weight matrix applied to a batch
  of 16 inputs, has
  \[
    I = \frac{4096 \cdot 16 \cdot 4096}
             {4096 \cdot 4096 + 16 \cdot 4096 + 4096 \cdot 16}
      = \frac{16 \cdot 4096}{4096 + 32}
      \approx \SI{15.9}{\flop\per\byte} ,
  \]
  far below the ratio: it is memory-limited, and a faster arithmetic
  unit would not speed it up.  The small dimension $N$ caps the
  intensity, since $I < N$ whenever $N \le M, K$.\tip{For a skinny
  GEMM with $N \ll M, K$, $I \approx N$ \si{\flop\per\byte} in FP16 (and $2N$ in
  FP8): read the intensity off the smallest dimension.}
\end{example}

The GEMM guide adds two effects that \cref{eq:l12-ai} does not
capture, because a GPU computes $C$ in tiles, each tile assigned to one
thread block.  The first, \term{tile quantization}[タイル量子化], \enquote{occurs
when matrix dimensions are not divisible by the thread block tile size}
\cite[§3.1]{nvidia2026gemmguide}: the last tile of a row or column is
only partly useful but costs as much as a full one.  The second,
\term{wave quantization}[ウェーブ量子化], arises because \enquote{the
total number of tiles is quantized to the number of multiprocessors on
the GPU} \cite[§3.2]{nvidia2026gemmguide}: tiles run in waves of one
tile per SM, and a last wave with few tiles leaves most SMs
idle.\margin{The guide's example: on an A100 with 108 SMs,
$9 \times 12 = 108$ tiles make one full wave, while $9 \times 13 = 117$
tiles add a second wave of only 9 tiles
\cite[§3.2]{nvidia2026gemmguide}.}  Tensor cores are also most efficient
when the GEMM dimensions are multiples of 8 elements for FP16 (on A100,
multiples of 64) \cite[§2.2]{nvidia2026gemmguide}.

\subsection{SIMD lanes versus a matrix unit}

On ordinary SIMT lanes, a matrix product is a loop of scalar
multiply-adds.  The kernel of \cref{lst:l12-matmul4} computes a
$4 \times 4$ product with one thread per element of $C$; each thread
runs the loop over $k$ sequentially and keeps its running sum in a
private register.

\begin{lstlisting}[language=C++, numbers=none, caption={A $4 \times 4$ matrix product on SIMT lanes, one thread per output element.}, label={lst:l12-matmul4}]
__global__ void matmul4(const float *A,
                        const float *B, float *C)
{
  const int N = 4;
  int row = threadIdx.y, col = threadIdx.x;
  float sum = 0.0f;           // private register
  for (int i = 0; i < N; ++i)
    sum += A[row * N + i] * B[i * N + col];
  C[row * N + col] = sum;
}
// launch: matmul4<<<1, dim3(4, 4)>>>(dA, dB, dC);
\end{lstlisting}

Because a CUDA thread's registers are private, one thread must perform
the whole accumulation of an output element.  Fewer threads (four, say,
one per row) simply loop longer.  Spreading one dot product over
several threads instead, for instance 16 threads that each compute one
product $a_{ik} b_{kj}$, requires a reduction of the partial sums
through shared memory or warp shuffles (Lesson~5), which adds
instructions and synchronisation.\margin{The reduction patterns of
Lesson~5 (tree reduction in shared memory, \texttt{\_\_shfl\_xor\_sync})
are exactly this overhead.}

A SIMD datapath can at best compute one element per operation: four
multipliers take a row of $A$ and a column of $B$, an adder tree sums
the four products, and the result is accumulated
(\cref{fig:l12-dot-unit}).  Row~1 of $C$ alone takes four of them,
\[
  c_{1j} = \sum_{k=1}^{4} a_{1k} b_{kj} , \qquad j = 1, \ldots, 4 ,
\]
each the dot product of row~1 of $A$ with one column of $B$, and the
product of two $4 \times 4$ matrices needs 16 such operations, one per
output element.

\begin{figure}[tb]
  \centering
  % One SIMD dot-product step for a 4x4 matrix product: four multipliers
% take row 1 of A and column 1 of B, an adder tree sums the four products,
% and the result is added to the accumulator c11.  A 4x4 product needs 16
% such steps, one per output element.
% Redraws the outline figure m12l1_simd_matrix.
% Usage: % One SIMD dot-product step for a 4x4 matrix product: four multipliers
% take row 1 of A and column 1 of B, an adder tree sums the four products,
% and the result is added to the accumulator c11.  A 4x4 product needs 16
% such steps, one per output element.
% Redraws the outline figure m12l1_simd_matrix.
% Usage: % One SIMD dot-product step for a 4x4 matrix product: four multipliers
% take row 1 of A and column 1 of B, an adder tree sums the four products,
% and the result is added to the accumulator c11.  A 4x4 product needs 16
% such steps, one per output element.
% Redraws the outline figure m12l1_simd_matrix.
% Usage: \input{figures/l12-dot-unit}
% Width: body (120 mm); the drawing is about 85 mm wide.
\begin{tikzpicture}[font=\scriptsize\sffamily,
    op/.style={draw, circle, fill=ln-accent!12, minimum size=6mm,
      inner sep=0pt},
    >={Stealth[length=1.3mm]}]
  \foreach \k [count=\n from 0] in {1,2,3,4} {
    \node[op] (m\k) at (\n*20mm,0) {$\times$};
    \node (a\k) at (\n*20mm-4.5mm,9mm) {$a_{1\k}$};
    \node (b\k) at (\n*20mm+4.5mm,9mm) {$b_{\k 1}$};
    \draw[->] (a\k) -- (m\k);
    \draw[->] (b\k) -- (m\k);
  }
  \node[op] (s1) at (10mm,-11mm) {$+$};
  \node[op] (s2) at (50mm,-11mm) {$+$};
  \node[op] (s3) at (30mm,-21mm) {$+$};
  \node[op, fill=ln-tip!20] (acc) at (30mm,-31mm) {$+$};
  \draw[->] (m1) -- (s1); \draw[->] (m2) -- (s1);
  \draw[->] (m3) -- (s2); \draw[->] (m4) -- (s2);
  \draw[->] (s1) -- (s3); \draw[->] (s2) -- (s3);
  \draw[->] (s3) -- (acc);
  \node (cin) at (5mm,-31mm) {$c_{11}$};
  \draw[->] (cin) -- (acc);
  \node (cout) at (55mm,-31mm) {$c_{11}$};
  \draw[->] (acc) -- (cout);
  \node[anchor=west, align=left, text=ln-muted] at (64mm,-16mm)
    {\iflangja{乗算器4個\\加算木\\アキュムレータ}{4 multipliers\\adder tree\\accumulate}};
\end{tikzpicture}

% Width: body (120 mm); the drawing is about 85 mm wide.
\begin{tikzpicture}[font=\scriptsize\sffamily,
    op/.style={draw, circle, fill=ln-accent!12, minimum size=6mm,
      inner sep=0pt},
    >={Stealth[length=1.3mm]}]
  \foreach \k [count=\n from 0] in {1,2,3,4} {
    \node[op] (m\k) at (\n*20mm,0) {$\times$};
    \node (a\k) at (\n*20mm-4.5mm,9mm) {$a_{1\k}$};
    \node (b\k) at (\n*20mm+4.5mm,9mm) {$b_{\k 1}$};
    \draw[->] (a\k) -- (m\k);
    \draw[->] (b\k) -- (m\k);
  }
  \node[op] (s1) at (10mm,-11mm) {$+$};
  \node[op] (s2) at (50mm,-11mm) {$+$};
  \node[op] (s3) at (30mm,-21mm) {$+$};
  \node[op, fill=ln-tip!20] (acc) at (30mm,-31mm) {$+$};
  \draw[->] (m1) -- (s1); \draw[->] (m2) -- (s1);
  \draw[->] (m3) -- (s2); \draw[->] (m4) -- (s2);
  \draw[->] (s1) -- (s3); \draw[->] (s2) -- (s3);
  \draw[->] (s3) -- (acc);
  \node (cin) at (5mm,-31mm) {$c_{11}$};
  \draw[->] (cin) -- (acc);
  \node (cout) at (55mm,-31mm) {$c_{11}$};
  \draw[->] (acc) -- (cout);
  \node[anchor=west, align=left, text=ln-muted] at (64mm,-16mm)
    {\iflangja{乗算器4個\\加算木\\アキュムレータ}{4 multipliers\\adder tree\\accumulate}};
\end{tikzpicture}

% Width: body (120 mm); the drawing is about 85 mm wide.
\begin{tikzpicture}[font=\scriptsize\sffamily,
    op/.style={draw, circle, fill=ln-accent!12, minimum size=6mm,
      inner sep=0pt},
    >={Stealth[length=1.3mm]}]
  \foreach \k [count=\n from 0] in {1,2,3,4} {
    \node[op] (m\k) at (\n*20mm,0) {$\times$};
    \node (a\k) at (\n*20mm-4.5mm,9mm) {$a_{1\k}$};
    \node (b\k) at (\n*20mm+4.5mm,9mm) {$b_{\k 1}$};
    \draw[->] (a\k) -- (m\k);
    \draw[->] (b\k) -- (m\k);
  }
  \node[op] (s1) at (10mm,-11mm) {$+$};
  \node[op] (s2) at (50mm,-11mm) {$+$};
  \node[op] (s3) at (30mm,-21mm) {$+$};
  \node[op, fill=ln-tip!20] (acc) at (30mm,-31mm) {$+$};
  \draw[->] (m1) -- (s1); \draw[->] (m2) -- (s1);
  \draw[->] (m3) -- (s2); \draw[->] (m4) -- (s2);
  \draw[->] (s1) -- (s3); \draw[->] (s2) -- (s3);
  \draw[->] (s3) -- (acc);
  \node (cin) at (5mm,-31mm) {$c_{11}$};
  \draw[->] (cin) -- (acc);
  \node (cout) at (55mm,-31mm) {$c_{11}$};
  \draw[->] (acc) -- (cout);
  \node[anchor=west, align=left, text=ln-muted] at (64mm,-16mm)
    {\iflangja{乗算器4個\\加算木\\アキュムレータ}{4 multipliers\\adder tree\\accumulate}};
\end{tikzpicture}

  \caption{One SIMD dot-product step: four multipliers, an adder tree and
    the accumulation compute $c_{11} \mathrel{+}= \sum_k a_{1k} b_{k1}$.
    A $4 \times 4$ product needs 16 such steps.}
  \label{fig:l12-dot-unit}
\end{figure}

\begin{definition}[MAC, MMA, tensor core instruction]
  A \term{multiply-accumulate (MAC)}[積和演算（MAC）] adds a product to
  a running sum, $c \leftarrow c + a b$; done with a single rounding it is
  the FMA.\caution{Lecture vs ISA: \texttt{HMMA.1688} names the warp-level shape
$m16n8k8$ (a $16 \times 8$ result, $K = 8$), the PTX \texttt{mma}
shape \texttt{.m16n8k8} \cite{nvidia2026ptx}, not a $4 \times 4$
product.  The $4 \times 4 \times 4$ figure is the per-clock datapath of
one Volta tensor core \cite{nvidia2020ampere}.}  A
  \term{matrix multiply-accumulate (MMA)}[行列積和演算（MMA）] applies
  the same idea to whole matrices, $D \leftarrow A B + C$, as one
  operation.  The tensor core (Lesson~1) is NVIDIA's MMA unit; the GPU
  has hundreds of them.  In NVIDIA's machine code (SASS) the half-precision MMA
  instruction is \term{HMMA}[HMMA], for instance
  \texttt{HMMA.1688.F16}: one instruction does the work that the SIMT
  loop does with many.
\end{definition}


The 32 threads of a warp issue the HMMA together, each supplying a
fragment of the operands from its registers, so the tensor core
sidesteps the private-register problem above.\margin{CUDA C++ reaches
the tensor cores through the warp-level matrix functions
(\texttt{wmma::load\_matrix\_sync}, \texttt{wmma::mma\_sync}): CPG §7.24
Warp Matrix Functions.  PTX and SASS are the subject of Lesson~10.}
Tensor-core generations differ mainly in the size of the product they
complete per clock and in the formats they accept
(\cref{tab:l12-tensor-cores}).  The A100 tensor core \enquote{can execute
256 FP16 FMA operations per clock}, an $8 \times 4 \times 8$
mixed-precision product, and accepts TF32, BF16 and FP64 inputs as
well as INT8, INT4 and binary (1-bit) \cite{nvidia2020ampere}; the
integer and binary modes date from Turing, the generation between the
two; sparsity support doubles the peak
again (\cref{sec:l12-tensor-core}).

\begin{table}[htbp]
  \centering
  \caption{Tensor cores of three NVIDIA generations (AWP
    \cite{nvidia2020ampere}; HWP V1.02, SXM5, preliminary
    \cite{nvidia2022hopper}).  Peak rates are dense / with 2:4
    sparsity.}
  \label{tab:l12-tensor-cores}
  \setlength{\tabcolsep}{4pt}
  \begin{tabular}{>{\raggedright\arraybackslash}p{46mm}lll}
    \toprule
    & V100 & A100 & H100 \\
    \midrule
    Tensor cores per SM & 8 & 4 & 4 \\
    FP16 FMA per clock per core & 64 & 256 & 512 \\
    Product per core per clock & $4 \times 4 \times 4$ & $8 \times 4 \times 8$ & --- \\
    \addlinespace
    FP16 tensor TFLOP/s & 125 & 312 / 624 & 1000 / 2000 \\
    TF32 tensor TFLOP/s & --- & 156 / 312 & 500 / 1000 \\
    FP8 tensor TFLOP/s & --- & --- & 2000 / 4000 \\
    FP32 (non-tensor) TFLOP/s & 15.7 & 19.5 & 60 \\
    \bottomrule
  \end{tabular}
\end{table}

\begin{figure}[tb]
  \centering
  % A 4x4 output-stationary systolic array at the start of a 4x4x4 matrix
% product C = AB.  Rows of A enter from the left and columns of B from the
% top, each row and column delayed (skewed) by one cycle relative to the
% previous one, so that a_ik and b_kj meet in PE(i,j) at cycle i+j+k-2.
% Each PE multiplies its two inputs, adds the product to its stationary
% partial sum c_ij, and passes a to the right and b downwards.
% Redraws the outline figure m12l1_matrix2.
% Usage: % A 4x4 output-stationary systolic array at the start of a 4x4x4 matrix
% product C = AB.  Rows of A enter from the left and columns of B from the
% top, each row and column delayed (skewed) by one cycle relative to the
% previous one, so that a_ik and b_kj meet in PE(i,j) at cycle i+j+k-2.
% Each PE multiplies its two inputs, adds the product to its stationary
% partial sum c_ij, and passes a to the right and b downwards.
% Redraws the outline figure m12l1_matrix2.
% Usage: % A 4x4 output-stationary systolic array at the start of a 4x4x4 matrix
% product C = AB.  Rows of A enter from the left and columns of B from the
% top, each row and column delayed (skewed) by one cycle relative to the
% previous one, so that a_ik and b_kj meet in PE(i,j) at cycle i+j+k-2.
% Each PE multiplies its two inputs, adds the product to its stationary
% partial sum c_ij, and passes a to the right and b downwards.
% Redraws the outline figure m12l1_matrix2.
% Usage: \input{figures/l12-systolic}
% Width: body (120 mm); the drawing is about 100 mm wide.
\begin{tikzpicture}[font=\scriptsize\sffamily,
    pe/.style={draw, fill=ln-accent!10, minimum size=9mm, inner sep=0pt,
      align=center},
    q/.style={inner sep=0pt, minimum width=5mm, minimum height=4mm},
    >={Stealth[length=1.3mm]}]
  \node[pe] (p11) at (13mm,-13mm) {PE\\$c_{11}$};
  \node[pe] (p12) at (26mm,-13mm) {PE\\$c_{12}$};
  \node[pe] (p13) at (39mm,-13mm) {PE\\$c_{13}$};
  \node[pe] (p14) at (52mm,-13mm) {PE\\$c_{14}$};
  \node[pe] (p21) at (13mm,-26mm) {PE\\$c_{21}$};
  \node[pe] (p22) at (26mm,-26mm) {PE\\$c_{22}$};
  \node[pe] (p23) at (39mm,-26mm) {PE\\$c_{23}$};
  \node[pe] (p24) at (52mm,-26mm) {PE\\$c_{24}$};
  \node[pe] (p31) at (13mm,-39mm) {PE\\$c_{31}$};
  \node[pe] (p32) at (26mm,-39mm) {PE\\$c_{32}$};
  \node[pe] (p33) at (39mm,-39mm) {PE\\$c_{33}$};
  \node[pe] (p34) at (52mm,-39mm) {PE\\$c_{34}$};
  \node[pe] (p41) at (13mm,-52mm) {PE\\$c_{41}$};
  \node[pe] (p42) at (26mm,-52mm) {PE\\$c_{42}$};
  \node[pe] (p43) at (39mm,-52mm) {PE\\$c_{43}$};
  \node[pe] (p44) at (52mm,-52mm) {PE\\$c_{44}$};
  \draw[->] (p11) -- (p12);
  \draw[->] (p12) -- (p13);
  \draw[->] (p13) -- (p14);
  \draw[->] (p21) -- (p22);
  \draw[->] (p22) -- (p23);
  \draw[->] (p23) -- (p24);
  \draw[->] (p31) -- (p32);
  \draw[->] (p32) -- (p33);
  \draw[->] (p33) -- (p34);
  \draw[->] (p41) -- (p42);
  \draw[->] (p42) -- (p43);
  \draw[->] (p43) -- (p44);
  \draw[->] (p11) -- (p21);
  \draw[->] (p21) -- (p31);
  \draw[->] (p31) -- (p41);
  \draw[->] (p12) -- (p22);
  \draw[->] (p22) -- (p32);
  \draw[->] (p32) -- (p42);
  \draw[->] (p13) -- (p23);
  \draw[->] (p23) -- (p33);
  \draw[->] (p33) -- (p43);
  \draw[->] (p14) -- (p24);
  \draw[->] (p24) -- (p34);
  \draw[->] (p34) -- (p44);
  \draw[->] (6.3mm,-13mm) -- (p11);
  \node[q] at (3.8mm,-13mm) {$a_{11}$};
  \node[q] at (-1.4mm,-13mm) {$a_{12}$};
  \node[q] at (-6.6mm,-13mm) {$a_{13}$};
  \node[q] at (-11.8mm,-13mm) {$a_{14}$};
  \node[q] at (-17.0mm,-13mm) {$\cdot$};
  \node[q] at (-22.2mm,-13mm) {$\cdot$};
  \node[q] at (-27.4mm,-13mm) {$\cdot$};
  \draw[->] (6.3mm,-26mm) -- (p21);
  \node[q] at (3.8mm,-26mm) {$\cdot$};
  \node[q] at (-1.4mm,-26mm) {$a_{21}$};
  \node[q] at (-6.6mm,-26mm) {$a_{22}$};
  \node[q] at (-11.8mm,-26mm) {$a_{23}$};
  \node[q] at (-17.0mm,-26mm) {$a_{24}$};
  \node[q] at (-22.2mm,-26mm) {$\cdot$};
  \node[q] at (-27.4mm,-26mm) {$\cdot$};
  \draw[->] (6.3mm,-39mm) -- (p31);
  \node[q] at (3.8mm,-39mm) {$\cdot$};
  \node[q] at (-1.4mm,-39mm) {$\cdot$};
  \node[q] at (-6.6mm,-39mm) {$a_{31}$};
  \node[q] at (-11.8mm,-39mm) {$a_{32}$};
  \node[q] at (-17.0mm,-39mm) {$a_{33}$};
  \node[q] at (-22.2mm,-39mm) {$a_{34}$};
  \node[q] at (-27.4mm,-39mm) {$\cdot$};
  \draw[->] (6.3mm,-52mm) -- (p41);
  \node[q] at (3.8mm,-52mm) {$\cdot$};
  \node[q] at (-1.4mm,-52mm) {$\cdot$};
  \node[q] at (-6.6mm,-52mm) {$\cdot$};
  \node[q] at (-11.8mm,-52mm) {$a_{41}$};
  \node[q] at (-17.0mm,-52mm) {$a_{42}$};
  \node[q] at (-22.2mm,-52mm) {$a_{43}$};
  \node[q] at (-27.4mm,-52mm) {$a_{44}$};
  \draw[->] (13mm,-6.3mm) -- (p11);
  \node[q] at (13mm,-4.3mm) {$b_{11}$};
  \node[q] at (13mm,-0.1mm) {$b_{21}$};
  \node[q] at (13mm,4.1mm) {$b_{31}$};
  \node[q] at (13mm,8.3mm) {$b_{41}$};
  \node[q] at (13mm,12.5mm) {$\cdot$};
  \node[q] at (13mm,16.7mm) {$\cdot$};
  \node[q] at (13mm,20.9mm) {$\cdot$};
  \draw[->] (26mm,-6.3mm) -- (p12);
  \node[q] at (26mm,-4.3mm) {$\cdot$};
  \node[q] at (26mm,-0.1mm) {$b_{12}$};
  \node[q] at (26mm,4.1mm) {$b_{22}$};
  \node[q] at (26mm,8.3mm) {$b_{32}$};
  \node[q] at (26mm,12.5mm) {$b_{42}$};
  \node[q] at (26mm,16.7mm) {$\cdot$};
  \node[q] at (26mm,20.9mm) {$\cdot$};
  \draw[->] (39mm,-6.3mm) -- (p13);
  \node[q] at (39mm,-4.3mm) {$\cdot$};
  \node[q] at (39mm,-0.1mm) {$\cdot$};
  \node[q] at (39mm,4.1mm) {$b_{13}$};
  \node[q] at (39mm,8.3mm) {$b_{23}$};
  \node[q] at (39mm,12.5mm) {$b_{33}$};
  \node[q] at (39mm,16.7mm) {$b_{43}$};
  \node[q] at (39mm,20.9mm) {$\cdot$};
  \draw[->] (52mm,-6.3mm) -- (p14);
  \node[q] at (52mm,-4.3mm) {$\cdot$};
  \node[q] at (52mm,-0.1mm) {$\cdot$};
  \node[q] at (52mm,4.1mm) {$\cdot$};
  \node[q] at (52mm,8.3mm) {$b_{14}$};
  \node[q] at (52mm,12.5mm) {$b_{24}$};
  \node[q] at (52mm,16.7mm) {$b_{34}$};
  \node[q] at (52mm,20.9mm) {$b_{44}$};
  % annotations
  \node[anchor=north, align=center] at (-6mm,-59mm)
    {\iflangja{$A$ の行（行ごとに1サイクル遅延）}{rows of $A$, each delayed one cycle}};
  \node[anchor=west, align=left] at (58mm,20mm)
    {\iflangja{$B$ の列\\（列ごとに1サイクル遅延）}{columns of $B$,\\each delayed one cycle}};
  \node[anchor=west, align=left] at (58mm,-32mm)
    {\iflangja{各 PE：\\$c_{ij} \mathrel{+}= a \cdot b$\\$a$ を右へ，$b$ を下へ}{each PE:\\$c_{ij} \mathrel{+}= a \cdot b$\\pass $a$ right, $b$ down}};
\end{tikzpicture}

% Width: body (120 mm); the drawing is about 100 mm wide.
\begin{tikzpicture}[font=\scriptsize\sffamily,
    pe/.style={draw, fill=ln-accent!10, minimum size=9mm, inner sep=0pt,
      align=center},
    q/.style={inner sep=0pt, minimum width=5mm, minimum height=4mm},
    >={Stealth[length=1.3mm]}]
  \node[pe] (p11) at (13mm,-13mm) {PE\\$c_{11}$};
  \node[pe] (p12) at (26mm,-13mm) {PE\\$c_{12}$};
  \node[pe] (p13) at (39mm,-13mm) {PE\\$c_{13}$};
  \node[pe] (p14) at (52mm,-13mm) {PE\\$c_{14}$};
  \node[pe] (p21) at (13mm,-26mm) {PE\\$c_{21}$};
  \node[pe] (p22) at (26mm,-26mm) {PE\\$c_{22}$};
  \node[pe] (p23) at (39mm,-26mm) {PE\\$c_{23}$};
  \node[pe] (p24) at (52mm,-26mm) {PE\\$c_{24}$};
  \node[pe] (p31) at (13mm,-39mm) {PE\\$c_{31}$};
  \node[pe] (p32) at (26mm,-39mm) {PE\\$c_{32}$};
  \node[pe] (p33) at (39mm,-39mm) {PE\\$c_{33}$};
  \node[pe] (p34) at (52mm,-39mm) {PE\\$c_{34}$};
  \node[pe] (p41) at (13mm,-52mm) {PE\\$c_{41}$};
  \node[pe] (p42) at (26mm,-52mm) {PE\\$c_{42}$};
  \node[pe] (p43) at (39mm,-52mm) {PE\\$c_{43}$};
  \node[pe] (p44) at (52mm,-52mm) {PE\\$c_{44}$};
  \draw[->] (p11) -- (p12);
  \draw[->] (p12) -- (p13);
  \draw[->] (p13) -- (p14);
  \draw[->] (p21) -- (p22);
  \draw[->] (p22) -- (p23);
  \draw[->] (p23) -- (p24);
  \draw[->] (p31) -- (p32);
  \draw[->] (p32) -- (p33);
  \draw[->] (p33) -- (p34);
  \draw[->] (p41) -- (p42);
  \draw[->] (p42) -- (p43);
  \draw[->] (p43) -- (p44);
  \draw[->] (p11) -- (p21);
  \draw[->] (p21) -- (p31);
  \draw[->] (p31) -- (p41);
  \draw[->] (p12) -- (p22);
  \draw[->] (p22) -- (p32);
  \draw[->] (p32) -- (p42);
  \draw[->] (p13) -- (p23);
  \draw[->] (p23) -- (p33);
  \draw[->] (p33) -- (p43);
  \draw[->] (p14) -- (p24);
  \draw[->] (p24) -- (p34);
  \draw[->] (p34) -- (p44);
  \draw[->] (6.3mm,-13mm) -- (p11);
  \node[q] at (3.8mm,-13mm) {$a_{11}$};
  \node[q] at (-1.4mm,-13mm) {$a_{12}$};
  \node[q] at (-6.6mm,-13mm) {$a_{13}$};
  \node[q] at (-11.8mm,-13mm) {$a_{14}$};
  \node[q] at (-17.0mm,-13mm) {$\cdot$};
  \node[q] at (-22.2mm,-13mm) {$\cdot$};
  \node[q] at (-27.4mm,-13mm) {$\cdot$};
  \draw[->] (6.3mm,-26mm) -- (p21);
  \node[q] at (3.8mm,-26mm) {$\cdot$};
  \node[q] at (-1.4mm,-26mm) {$a_{21}$};
  \node[q] at (-6.6mm,-26mm) {$a_{22}$};
  \node[q] at (-11.8mm,-26mm) {$a_{23}$};
  \node[q] at (-17.0mm,-26mm) {$a_{24}$};
  \node[q] at (-22.2mm,-26mm) {$\cdot$};
  \node[q] at (-27.4mm,-26mm) {$\cdot$};
  \draw[->] (6.3mm,-39mm) -- (p31);
  \node[q] at (3.8mm,-39mm) {$\cdot$};
  \node[q] at (-1.4mm,-39mm) {$\cdot$};
  \node[q] at (-6.6mm,-39mm) {$a_{31}$};
  \node[q] at (-11.8mm,-39mm) {$a_{32}$};
  \node[q] at (-17.0mm,-39mm) {$a_{33}$};
  \node[q] at (-22.2mm,-39mm) {$a_{34}$};
  \node[q] at (-27.4mm,-39mm) {$\cdot$};
  \draw[->] (6.3mm,-52mm) -- (p41);
  \node[q] at (3.8mm,-52mm) {$\cdot$};
  \node[q] at (-1.4mm,-52mm) {$\cdot$};
  \node[q] at (-6.6mm,-52mm) {$\cdot$};
  \node[q] at (-11.8mm,-52mm) {$a_{41}$};
  \node[q] at (-17.0mm,-52mm) {$a_{42}$};
  \node[q] at (-22.2mm,-52mm) {$a_{43}$};
  \node[q] at (-27.4mm,-52mm) {$a_{44}$};
  \draw[->] (13mm,-6.3mm) -- (p11);
  \node[q] at (13mm,-4.3mm) {$b_{11}$};
  \node[q] at (13mm,-0.1mm) {$b_{21}$};
  \node[q] at (13mm,4.1mm) {$b_{31}$};
  \node[q] at (13mm,8.3mm) {$b_{41}$};
  \node[q] at (13mm,12.5mm) {$\cdot$};
  \node[q] at (13mm,16.7mm) {$\cdot$};
  \node[q] at (13mm,20.9mm) {$\cdot$};
  \draw[->] (26mm,-6.3mm) -- (p12);
  \node[q] at (26mm,-4.3mm) {$\cdot$};
  \node[q] at (26mm,-0.1mm) {$b_{12}$};
  \node[q] at (26mm,4.1mm) {$b_{22}$};
  \node[q] at (26mm,8.3mm) {$b_{32}$};
  \node[q] at (26mm,12.5mm) {$b_{42}$};
  \node[q] at (26mm,16.7mm) {$\cdot$};
  \node[q] at (26mm,20.9mm) {$\cdot$};
  \draw[->] (39mm,-6.3mm) -- (p13);
  \node[q] at (39mm,-4.3mm) {$\cdot$};
  \node[q] at (39mm,-0.1mm) {$\cdot$};
  \node[q] at (39mm,4.1mm) {$b_{13}$};
  \node[q] at (39mm,8.3mm) {$b_{23}$};
  \node[q] at (39mm,12.5mm) {$b_{33}$};
  \node[q] at (39mm,16.7mm) {$b_{43}$};
  \node[q] at (39mm,20.9mm) {$\cdot$};
  \draw[->] (52mm,-6.3mm) -- (p14);
  \node[q] at (52mm,-4.3mm) {$\cdot$};
  \node[q] at (52mm,-0.1mm) {$\cdot$};
  \node[q] at (52mm,4.1mm) {$\cdot$};
  \node[q] at (52mm,8.3mm) {$b_{14}$};
  \node[q] at (52mm,12.5mm) {$b_{24}$};
  \node[q] at (52mm,16.7mm) {$b_{34}$};
  \node[q] at (52mm,20.9mm) {$b_{44}$};
  % annotations
  \node[anchor=north, align=center] at (-6mm,-59mm)
    {\iflangja{$A$ の行（行ごとに1サイクル遅延）}{rows of $A$, each delayed one cycle}};
  \node[anchor=west, align=left] at (58mm,20mm)
    {\iflangja{$B$ の列\\（列ごとに1サイクル遅延）}{columns of $B$,\\each delayed one cycle}};
  \node[anchor=west, align=left] at (58mm,-32mm)
    {\iflangja{各 PE：\\$c_{ij} \mathrel{+}= a \cdot b$\\$a$ を右へ，$b$ を下へ}{each PE:\\$c_{ij} \mathrel{+}= a \cdot b$\\pass $a$ right, $b$ down}};
\end{tikzpicture}

% Width: body (120 mm); the drawing is about 100 mm wide.
\begin{tikzpicture}[font=\scriptsize\sffamily,
    pe/.style={draw, fill=ln-accent!10, minimum size=9mm, inner sep=0pt,
      align=center},
    q/.style={inner sep=0pt, minimum width=5mm, minimum height=4mm},
    >={Stealth[length=1.3mm]}]
  \node[pe] (p11) at (13mm,-13mm) {PE\\$c_{11}$};
  \node[pe] (p12) at (26mm,-13mm) {PE\\$c_{12}$};
  \node[pe] (p13) at (39mm,-13mm) {PE\\$c_{13}$};
  \node[pe] (p14) at (52mm,-13mm) {PE\\$c_{14}$};
  \node[pe] (p21) at (13mm,-26mm) {PE\\$c_{21}$};
  \node[pe] (p22) at (26mm,-26mm) {PE\\$c_{22}$};
  \node[pe] (p23) at (39mm,-26mm) {PE\\$c_{23}$};
  \node[pe] (p24) at (52mm,-26mm) {PE\\$c_{24}$};
  \node[pe] (p31) at (13mm,-39mm) {PE\\$c_{31}$};
  \node[pe] (p32) at (26mm,-39mm) {PE\\$c_{32}$};
  \node[pe] (p33) at (39mm,-39mm) {PE\\$c_{33}$};
  \node[pe] (p34) at (52mm,-39mm) {PE\\$c_{34}$};
  \node[pe] (p41) at (13mm,-52mm) {PE\\$c_{41}$};
  \node[pe] (p42) at (26mm,-52mm) {PE\\$c_{42}$};
  \node[pe] (p43) at (39mm,-52mm) {PE\\$c_{43}$};
  \node[pe] (p44) at (52mm,-52mm) {PE\\$c_{44}$};
  \draw[->] (p11) -- (p12);
  \draw[->] (p12) -- (p13);
  \draw[->] (p13) -- (p14);
  \draw[->] (p21) -- (p22);
  \draw[->] (p22) -- (p23);
  \draw[->] (p23) -- (p24);
  \draw[->] (p31) -- (p32);
  \draw[->] (p32) -- (p33);
  \draw[->] (p33) -- (p34);
  \draw[->] (p41) -- (p42);
  \draw[->] (p42) -- (p43);
  \draw[->] (p43) -- (p44);
  \draw[->] (p11) -- (p21);
  \draw[->] (p21) -- (p31);
  \draw[->] (p31) -- (p41);
  \draw[->] (p12) -- (p22);
  \draw[->] (p22) -- (p32);
  \draw[->] (p32) -- (p42);
  \draw[->] (p13) -- (p23);
  \draw[->] (p23) -- (p33);
  \draw[->] (p33) -- (p43);
  \draw[->] (p14) -- (p24);
  \draw[->] (p24) -- (p34);
  \draw[->] (p34) -- (p44);
  \draw[->] (6.3mm,-13mm) -- (p11);
  \node[q] at (3.8mm,-13mm) {$a_{11}$};
  \node[q] at (-1.4mm,-13mm) {$a_{12}$};
  \node[q] at (-6.6mm,-13mm) {$a_{13}$};
  \node[q] at (-11.8mm,-13mm) {$a_{14}$};
  \node[q] at (-17.0mm,-13mm) {$\cdot$};
  \node[q] at (-22.2mm,-13mm) {$\cdot$};
  \node[q] at (-27.4mm,-13mm) {$\cdot$};
  \draw[->] (6.3mm,-26mm) -- (p21);
  \node[q] at (3.8mm,-26mm) {$\cdot$};
  \node[q] at (-1.4mm,-26mm) {$a_{21}$};
  \node[q] at (-6.6mm,-26mm) {$a_{22}$};
  \node[q] at (-11.8mm,-26mm) {$a_{23}$};
  \node[q] at (-17.0mm,-26mm) {$a_{24}$};
  \node[q] at (-22.2mm,-26mm) {$\cdot$};
  \node[q] at (-27.4mm,-26mm) {$\cdot$};
  \draw[->] (6.3mm,-39mm) -- (p31);
  \node[q] at (3.8mm,-39mm) {$\cdot$};
  \node[q] at (-1.4mm,-39mm) {$\cdot$};
  \node[q] at (-6.6mm,-39mm) {$a_{31}$};
  \node[q] at (-11.8mm,-39mm) {$a_{32}$};
  \node[q] at (-17.0mm,-39mm) {$a_{33}$};
  \node[q] at (-22.2mm,-39mm) {$a_{34}$};
  \node[q] at (-27.4mm,-39mm) {$\cdot$};
  \draw[->] (6.3mm,-52mm) -- (p41);
  \node[q] at (3.8mm,-52mm) {$\cdot$};
  \node[q] at (-1.4mm,-52mm) {$\cdot$};
  \node[q] at (-6.6mm,-52mm) {$\cdot$};
  \node[q] at (-11.8mm,-52mm) {$a_{41}$};
  \node[q] at (-17.0mm,-52mm) {$a_{42}$};
  \node[q] at (-22.2mm,-52mm) {$a_{43}$};
  \node[q] at (-27.4mm,-52mm) {$a_{44}$};
  \draw[->] (13mm,-6.3mm) -- (p11);
  \node[q] at (13mm,-4.3mm) {$b_{11}$};
  \node[q] at (13mm,-0.1mm) {$b_{21}$};
  \node[q] at (13mm,4.1mm) {$b_{31}$};
  \node[q] at (13mm,8.3mm) {$b_{41}$};
  \node[q] at (13mm,12.5mm) {$\cdot$};
  \node[q] at (13mm,16.7mm) {$\cdot$};
  \node[q] at (13mm,20.9mm) {$\cdot$};
  \draw[->] (26mm,-6.3mm) -- (p12);
  \node[q] at (26mm,-4.3mm) {$\cdot$};
  \node[q] at (26mm,-0.1mm) {$b_{12}$};
  \node[q] at (26mm,4.1mm) {$b_{22}$};
  \node[q] at (26mm,8.3mm) {$b_{32}$};
  \node[q] at (26mm,12.5mm) {$b_{42}$};
  \node[q] at (26mm,16.7mm) {$\cdot$};
  \node[q] at (26mm,20.9mm) {$\cdot$};
  \draw[->] (39mm,-6.3mm) -- (p13);
  \node[q] at (39mm,-4.3mm) {$\cdot$};
  \node[q] at (39mm,-0.1mm) {$\cdot$};
  \node[q] at (39mm,4.1mm) {$b_{13}$};
  \node[q] at (39mm,8.3mm) {$b_{23}$};
  \node[q] at (39mm,12.5mm) {$b_{33}$};
  \node[q] at (39mm,16.7mm) {$b_{43}$};
  \node[q] at (39mm,20.9mm) {$\cdot$};
  \draw[->] (52mm,-6.3mm) -- (p14);
  \node[q] at (52mm,-4.3mm) {$\cdot$};
  \node[q] at (52mm,-0.1mm) {$\cdot$};
  \node[q] at (52mm,4.1mm) {$\cdot$};
  \node[q] at (52mm,8.3mm) {$b_{14}$};
  \node[q] at (52mm,12.5mm) {$b_{24}$};
  \node[q] at (52mm,16.7mm) {$b_{34}$};
  \node[q] at (52mm,20.9mm) {$b_{44}$};
  % annotations
  \node[anchor=north, align=center] at (-6mm,-59mm)
    {\iflangja{$A$ の行（行ごとに1サイクル遅延）}{rows of $A$, each delayed one cycle}};
  \node[anchor=west, align=left] at (58mm,20mm)
    {\iflangja{$B$ の列\\（列ごとに1サイクル遅延）}{columns of $B$,\\each delayed one cycle}};
  \node[anchor=west, align=left] at (58mm,-32mm)
    {\iflangja{各 PE：\\$c_{ij} \mathrel{+}= a \cdot b$\\$a$ を右へ，$b$ を下へ}{each PE:\\$c_{ij} \mathrel{+}= a \cdot b$\\pass $a$ right, $b$ down}};
\end{tikzpicture}

  \caption{A $4 \times 4$ output-stationary systolic array at the start of
    $C = AB$.  The skew makes $a_{ik}$ and $b_{kj}$ meet in PE$(i,j)$ in
    cycle $i + j + k - 2$; dots mark empty slots.}
  \label{fig:l12-systolic}
\end{figure}

\subsection{Systolic arrays}

The traditional hardware for matrix multiplication is the systolic
array.

\begin{definition}[Systolic array]
  A \term{systolic array}[シストリックアレイ] is a grid of identical
  processing elements (PEs), each a MAC unit, connected only to their
  neighbours.  Data enter at the edges and move one PE per cycle; every
  PE computes on the values that pass through it and forwards
  them.\sidenote{Kung named the structure after the rhythmic,
  pulsed flow of data through the array.  Each value fetched from
  memory is used by many PEs before it leaves, which is what makes the
  structure attractive when memory bandwidth is the limit
  \cite{kung1982systolic}.  Google's first TPU is built around a
  $256 \times 256$ array, \enquote{a 65,536 8-bit MAC matrix multiply
  unit} with a peak of 92~TOPS \cite{jouppi2017tpu}.}
\end{definition}

In the \emph{output-stationary} arrangement of
\cref{fig:l12-systolic}, PE$(i,j)$ owns the output element $c_{ij}$.
The rows of $A$ enter from the left and the columns of $B$ from the
top, each row and column delayed by one cycle relative to the previous
one (the \emph{skew}).  In every cycle each PE multiplies the $a$ and
$b$ that arrive, adds the product to its stationary partial sum, and
passes $a$ to the right and $b$ downwards.  In cycle~1 PE$(1,1)$
computes $a_{11} b_{11}$; in cycle~2 $a_{11}$ has moved right to meet
$b_{12}$ in PE$(1,2)$, while PE$(1,1)$ adds $a_{12} b_{21}$; and so on.

With this skew, PE$(i,j)$ performs step $k$ in cycle $t = i + j + k - 2$.
For an $n \times n$ array the first PE starts in cycle~1, the last
PE$(n,n)$ starts in cycle $2n - 1$ and finishes in cycle $3n - 2$.
\Cref{tab:l12-systolic-trace} traces $n = 4$: the 64 MACs take 10
cycles, and no cycle has all 16 PEs busy.

\begin{table}[htbp]
  \centering
  \caption{Activity of the $4 \times 4$ array of \cref{fig:l12-systolic}
    for one $4 \times 4 \times 4$ product (our model:
    PE$(i,j)$ performs step $k$ in cycle $i + j + k - 2$).}
  \label{tab:l12-systolic-trace}
  \begin{tabular}{lcccccccccc}
    \toprule
    Cycle & 1 & 2 & 3 & 4 & 5 & 6 & 7 & 8 & 9 & 10 \\
    \midrule
    Busy PEs & 1 & 3 & 6 & 10 & 12 & 12 & 10 & 6 & 3 & 1 \\
    Outputs finished & -- & -- & -- & 1 & 2 & 3 & 4 & 3 & 2 & 1 \\
    \bottomrule
  \end{tabular}
\end{table}

\begin{definition}[Fill and drain]
  The \term{fill and drain}[充填と排出] of a systolic array are the
  cycles at the start of a computation, while the skewed inputs are
  still entering and only some PEs are busy, and at its end, while the
  last operands pass through and the PEs fall idle again.
\end{definition}

In the trace, the array does 64 MACs in 10 cycles on 16 PEs, a
utilisation of $64/160 = 40\,\%$; the rest is fill and drain.  When
many tiles are streamed through back to back, the next tile fills the
array while the previous one drains, and the utilisation approaches
100\,\%.\caution{The lecture describes the array as fully busy around
cycle~4 and finished in about seven cycles; it counts differently from
the model above.  Under our model, cycle~4 has 10 of 16 PEs busy,
PE$(4,4)$ starts in cycle $2n - 1 = 7$, and the product completes in
cycle $3n - 2 = 10$.}

\subsection{Transformer Engine and data movement}
\label{sec:l12-te}

Hopper's Transformer Engine (Lesson~1) accelerates the layers of
transformer models by scaling tensor data dynamically into the
representable range of a narrow format, which makes FP8 arithmetic
usable \cite{nvidia2022hopper}.\margin{\Cref{sec:l12-formats} explains
why a narrow format needs this scaling: FP8 E4M3 covers magnitudes only
from $2^{-9}$ to 448.}  Fast arithmetic is only half of the problem:
the units must be fed.  The Tensor Memory Accelerator (Lesson~1) moves
blocks of a tensor between global and shared memory in chunks and
generates the addresses itself.

\begin{definition}[Copy descriptor]
  A \term{copy descriptor}[コピー記述子] is a data structure that
  describes a multidimensional tensor and the blocks to be copied from it
  (base address, tensor sizes and strides, and the shape of the block), so
  that the TMA generates all addresses of a transfer from it.  The
  coordinates of a block are given with each copy instruction, so one
  descriptor serves all the tiles of a tensor.  In Hopper, \enquote{a
  single thread creates a copy descriptor} before issuing the transfer
  \cite{nvidia2022hopper}, instead of every thread computing its own
  addresses.
\end{definition}

\begin{keypoint}
  SIMD/SIMT lanes can accelerate matrix multiplication, dedicated
  matrix units accelerate it much further, and in both cases keeping
  the memory bandwidth fully used is critical: a GEMM whose intensity
  is below the ops:byte ratio gains nothing from a faster unit.
\end{keypoint}

\section{Floating-point formats}
\label{sec:l12-formats}

Scientific computing relies on IEEE single (32-bit) and double
(64-bit) precision; ML reduces the precision.\source{M12, floating-point
formats; PMPP App.~A; BPG §7.3 Numerical Accuracy and Precision; \cite{ieee754-2019,nvidia2020ampere,nvidia2022hopper}.}
This section explains why narrower formats pay, how a floating-point
number is encoded, and how ML-specific formats trade precision for
range.

\subsection{Why lower precision pays}

The number of floating-point operations of a GEMM depends on its shape,
not on its format.  The bytes moved depend on the format: halving the
input width halves the bytes and therefore doubles the arithmetic
intensity.  With FP8 inputs (one byte per element), \cref{eq:l12-ai}
becomes
\begin{equation}
  I_{\text{GEMM, FP8}} = \frac{2MNK}{MK + NK + MN} ,
  \label{eq:l12-ai-fp8}
\end{equation}
twice the FP16 value; the skinny GEMM of \cref{ex:l12-ai} rises from
15.9 to \SI{31.8}{\flop\per\byte}.\margin{\Cref{eq:l12-ai-fp8} counts $C$
at the input width as well; an FP32 output would add $2MN$ bytes.}

\begin{definition}[Quantization, mixed precision, accumulator]
  The representation of values with fewer bits than the format they were computed in, for instance FP32
  weights stored in FP8, is called \term{quantization}[量子化].  In
  \term{mixed precision}[混合精度] arithmetic the inputs of an operation
  are in a narrow format while the
  \term{accumulator}[アキュムレータ], the register or variable that
  holds the running sum, stays in a wider one, typically FP16 or FP8
  inputs with FP32 accumulation.
\end{definition}

Quantization brings three benefits: smaller storage, higher arithmetic
intensity, and, on recent hardware, higher arithmetic throughput,
because a unit of the same area completes more narrow operations per
cycle.  On older units only the first two applied: the operands shrank,
but the operations per cycle stayed the same.\margin{Hopper:
\enquote{FP8 halves data storage requirements and doubles throughput
compared to FP16 or BF16} \cite{nvidia2022hopper}; compare the FP16 and
FP8 rows of \cref{tab:l12-tensor-cores}.}\caution{The accumulator stays
wide for a reason: a long sum of narrow terms loses the small addends
(Lesson~5, lower precision).  Volta's tensor cores already took FP16 inputs with
\enquote{FP32 accumulation} \cite{nvidia2020ampere}.}

\subsection{Sign, exponent and fraction}

A floating-point number is the binary form of scientific notation.  In
decimal, $0.0012 = 1.2 \times 10^{-3}$ and
$0.0000012 = 1.2 \times 10^{-6}$ share the fraction $1.2$ and differ
only in the exponent; a binary format uses base~2 in the same way.

\begin{definition}[Floating-point format]
  A binary \term{floating-point format}[浮動小数点形式] stores a number
  in three fields: a \term{sign bit}[符号ビット] $s$, an
  \term{exponent}[指数部] field $e$ of $w$ bits, and a
  \term{mantissa (fraction)}[仮数部] field $f$ of $p$ bits.\margin{\enquote{Mantissa},
\enquote{fraction} and \enquote{significand} name the same field in
practice; IEEE 754 calls the stored field the trailing significand
\cite{ieee754-2019}.}  The total
  width is $1 + w + p$ bits.  Because the exponent must represent both
  negative and positive powers (values below and above~1) while the
  sign bit belongs to the value, the exponent is stored with an offset,
  the \term{exponent bias}[指数バイアス] $b = 2^{w-1} - 1$.  A
  \emph{normal} number has $0 < e < 2^w - 1$ and the value
  \begin{equation}
    x = (-1)^{s} \times \Bigl(1 + \frac{f}{2^{p}}\Bigr) \times 2^{e - b} ,
    \label{eq:l12-normal}
  \end{equation}
  where the leading 1 is implicit and not stored.
\end{definition}

More exponent bits cover a wider range; more fraction bits give more
precision.  \Cref{fig:l12-formats} encodes the same value in the six
formats of this chapter, and \cref{tab:l12-formats} gives their
parameters.  Formats need not have power-of-two widths:
TF32 is $1 + 8 + 10 = 19$ bits wide.

\begin{figure}[tb]
  \centering
  % Bit layouts of six floating-point formats, each encoding the value
% 0.15625 = 1.01 (binary) x 2^-3: FP32 (1/8/23), TF32 (1/8/10), BF16
% (1/8/7), FP16 (1/5/10), FP8 E5M2 (1/5/2) and FP8 E4M3 (1/4/3).  Sign,
% exponent and fraction fields are shaded differently; the biased exponent
% is 124 (bias 127), 12 (bias 15) or 4 (bias 7).
% Redraws the bfloat16 comparison inside the outline figure
% m12l2_floating_basics_diff (originally Wikipedia, CC BY-SA 4.0), redrawn
% with the formats of this chapter; the AMD fp24 and Pixar PXR24 rows are
% left out.
% Usage: % Bit layouts of six floating-point formats, each encoding the value
% 0.15625 = 1.01 (binary) x 2^-3: FP32 (1/8/23), TF32 (1/8/10), BF16
% (1/8/7), FP16 (1/5/10), FP8 E5M2 (1/5/2) and FP8 E4M3 (1/4/3).  Sign,
% exponent and fraction fields are shaded differently; the biased exponent
% is 124 (bias 127), 12 (bias 15) or 4 (bias 7).
% Redraws the bfloat16 comparison inside the outline figure
% m12l2_floating_basics_diff (originally Wikipedia, CC BY-SA 4.0), redrawn
% with the formats of this chapter; the AMD fp24 and Pixar PXR24 rows are
% left out.
% Usage: % Bit layouts of six floating-point formats, each encoding the value
% 0.15625 = 1.01 (binary) x 2^-3: FP32 (1/8/23), TF32 (1/8/10), BF16
% (1/8/7), FP16 (1/5/10), FP8 E5M2 (1/5/2) and FP8 E4M3 (1/4/3).  Sign,
% exponent and fraction fields are shaded differently; the biased exponent
% is 124 (bias 127), 12 (bias 15) or 4 (bias 7).
% Redraws the bfloat16 comparison inside the outline figure
% m12l2_floating_basics_diff (originally Wikipedia, CC BY-SA 4.0), redrawn
% with the formats of this chapter; the AMD fp24 and Pixar PXR24 rows are
% left out.
% Usage: \input{figures/l12-formats}
% Width: body (120 mm); the drawing is about 116 mm wide.
\begin{tikzpicture}[font=\scriptsize\sffamily,
    bit/.style={draw=ln-muted, minimum width=2.9mm, minimum height=4.2mm,
      inner sep=0pt, font=\scriptsize\ttfamily}]
  % FP32
  \node[anchor=east] at (22mm,0.0mm) {FP32\enspace\textcolor{ln-muted}{1/8/23}};
  \node[bit, fill=ln-caution!30] at (24.45mm,0.0mm) {0};
  \node[bit, fill=ln-accent!22] at (27.35mm,0.0mm) {0};
  \node[bit, fill=ln-accent!22] at (30.25mm,0.0mm) {1};
  \node[bit, fill=ln-accent!22] at (33.15mm,0.0mm) {1};
  \node[bit, fill=ln-accent!22] at (36.05mm,0.0mm) {1};
  \node[bit, fill=ln-accent!22] at (38.95mm,0.0mm) {1};
  \node[bit, fill=ln-accent!22] at (41.85mm,0.0mm) {1};
  \node[bit, fill=ln-accent!22] at (44.75mm,0.0mm) {0};
  \node[bit, fill=ln-accent!22] at (47.65mm,0.0mm) {0};
  \node[bit, fill=ln-tip!22] at (50.55mm,0.0mm) {0};
  \node[bit, fill=ln-tip!22] at (53.45mm,0.0mm) {1};
  \node[bit, fill=ln-tip!22] at (56.35mm,0.0mm) {0};
  \node[bit, fill=ln-tip!22] at (59.25mm,0.0mm) {0};
  \node[bit, fill=ln-tip!22] at (62.15mm,0.0mm) {0};
  \node[bit, fill=ln-tip!22] at (65.05mm,0.0mm) {0};
  \node[bit, fill=ln-tip!22] at (67.95mm,0.0mm) {0};
  \node[bit, fill=ln-tip!22] at (70.85mm,0.0mm) {0};
  \node[bit, fill=ln-tip!22] at (73.75mm,0.0mm) {0};
  \node[bit, fill=ln-tip!22] at (76.65mm,0.0mm) {0};
  \node[bit, fill=ln-tip!22] at (79.55mm,0.0mm) {0};
  \node[bit, fill=ln-tip!22] at (82.45mm,0.0mm) {0};
  \node[bit, fill=ln-tip!22] at (85.35mm,0.0mm) {0};
  \node[bit, fill=ln-tip!22] at (88.25mm,0.0mm) {0};
  \node[bit, fill=ln-tip!22] at (91.15mm,0.0mm) {0};
  \node[bit, fill=ln-tip!22] at (94.05mm,0.0mm) {0};
  \node[bit, fill=ln-tip!22] at (96.95mm,0.0mm) {0};
  \node[bit, fill=ln-tip!22] at (99.85mm,0.0mm) {0};
  \node[bit, fill=ln-tip!22] at (102.75mm,0.0mm) {0};
  \node[bit, fill=ln-tip!22] at (105.65mm,0.0mm) {0};
  \node[bit, fill=ln-tip!22] at (108.55mm,0.0mm) {0};
  \node[bit, fill=ln-tip!22] at (111.45mm,0.0mm) {0};
  \node[bit, fill=ln-tip!22] at (114.35mm,0.0mm) {0};
  % TF32
  \node[anchor=east] at (22mm,-6.0mm) {TF32\enspace\textcolor{ln-muted}{1/8/10}};
  \node[bit, fill=ln-caution!30] at (24.45mm,-6.0mm) {0};
  \node[bit, fill=ln-accent!22] at (27.35mm,-6.0mm) {0};
  \node[bit, fill=ln-accent!22] at (30.25mm,-6.0mm) {1};
  \node[bit, fill=ln-accent!22] at (33.15mm,-6.0mm) {1};
  \node[bit, fill=ln-accent!22] at (36.05mm,-6.0mm) {1};
  \node[bit, fill=ln-accent!22] at (38.95mm,-6.0mm) {1};
  \node[bit, fill=ln-accent!22] at (41.85mm,-6.0mm) {1};
  \node[bit, fill=ln-accent!22] at (44.75mm,-6.0mm) {0};
  \node[bit, fill=ln-accent!22] at (47.65mm,-6.0mm) {0};
  \node[bit, fill=ln-tip!22] at (50.55mm,-6.0mm) {0};
  \node[bit, fill=ln-tip!22] at (53.45mm,-6.0mm) {1};
  \node[bit, fill=ln-tip!22] at (56.35mm,-6.0mm) {0};
  \node[bit, fill=ln-tip!22] at (59.25mm,-6.0mm) {0};
  \node[bit, fill=ln-tip!22] at (62.15mm,-6.0mm) {0};
  \node[bit, fill=ln-tip!22] at (65.05mm,-6.0mm) {0};
  \node[bit, fill=ln-tip!22] at (67.95mm,-6.0mm) {0};
  \node[bit, fill=ln-tip!22] at (70.85mm,-6.0mm) {0};
  \node[bit, fill=ln-tip!22] at (73.75mm,-6.0mm) {0};
  \node[bit, fill=ln-tip!22] at (76.65mm,-6.0mm) {0};
  % BF16
  \node[anchor=east] at (22mm,-12.0mm) {BF16\enspace\textcolor{ln-muted}{1/8/7}};
  \node[bit, fill=ln-caution!30] at (24.45mm,-12.0mm) {0};
  \node[bit, fill=ln-accent!22] at (27.35mm,-12.0mm) {0};
  \node[bit, fill=ln-accent!22] at (30.25mm,-12.0mm) {1};
  \node[bit, fill=ln-accent!22] at (33.15mm,-12.0mm) {1};
  \node[bit, fill=ln-accent!22] at (36.05mm,-12.0mm) {1};
  \node[bit, fill=ln-accent!22] at (38.95mm,-12.0mm) {1};
  \node[bit, fill=ln-accent!22] at (41.85mm,-12.0mm) {1};
  \node[bit, fill=ln-accent!22] at (44.75mm,-12.0mm) {0};
  \node[bit, fill=ln-accent!22] at (47.65mm,-12.0mm) {0};
  \node[bit, fill=ln-tip!22] at (50.55mm,-12.0mm) {0};
  \node[bit, fill=ln-tip!22] at (53.45mm,-12.0mm) {1};
  \node[bit, fill=ln-tip!22] at (56.35mm,-12.0mm) {0};
  \node[bit, fill=ln-tip!22] at (59.25mm,-12.0mm) {0};
  \node[bit, fill=ln-tip!22] at (62.15mm,-12.0mm) {0};
  \node[bit, fill=ln-tip!22] at (65.05mm,-12.0mm) {0};
  \node[bit, fill=ln-tip!22] at (67.95mm,-12.0mm) {0};
  % FP16
  \node[anchor=east] at (22mm,-18.0mm) {FP16\enspace\textcolor{ln-muted}{1/5/10}};
  \node[bit, fill=ln-caution!30] at (24.45mm,-18.0mm) {0};
  \node[bit, fill=ln-accent!22] at (27.35mm,-18.0mm) {0};
  \node[bit, fill=ln-accent!22] at (30.25mm,-18.0mm) {1};
  \node[bit, fill=ln-accent!22] at (33.15mm,-18.0mm) {1};
  \node[bit, fill=ln-accent!22] at (36.05mm,-18.0mm) {0};
  \node[bit, fill=ln-accent!22] at (38.95mm,-18.0mm) {0};
  \node[bit, fill=ln-tip!22] at (41.85mm,-18.0mm) {0};
  \node[bit, fill=ln-tip!22] at (44.75mm,-18.0mm) {1};
  \node[bit, fill=ln-tip!22] at (47.65mm,-18.0mm) {0};
  \node[bit, fill=ln-tip!22] at (50.55mm,-18.0mm) {0};
  \node[bit, fill=ln-tip!22] at (53.45mm,-18.0mm) {0};
  \node[bit, fill=ln-tip!22] at (56.35mm,-18.0mm) {0};
  \node[bit, fill=ln-tip!22] at (59.25mm,-18.0mm) {0};
  \node[bit, fill=ln-tip!22] at (62.15mm,-18.0mm) {0};
  \node[bit, fill=ln-tip!22] at (65.05mm,-18.0mm) {0};
  \node[bit, fill=ln-tip!22] at (67.95mm,-18.0mm) {0};
  % FP8 E5M2
  \node[anchor=east] at (22mm,-24.0mm) {FP8 E5M2\enspace\textcolor{ln-muted}{1/5/2}};
  \node[bit, fill=ln-caution!30] at (24.45mm,-24.0mm) {0};
  \node[bit, fill=ln-accent!22] at (27.35mm,-24.0mm) {0};
  \node[bit, fill=ln-accent!22] at (30.25mm,-24.0mm) {1};
  \node[bit, fill=ln-accent!22] at (33.15mm,-24.0mm) {1};
  \node[bit, fill=ln-accent!22] at (36.05mm,-24.0mm) {0};
  \node[bit, fill=ln-accent!22] at (38.95mm,-24.0mm) {0};
  \node[bit, fill=ln-tip!22] at (41.85mm,-24.0mm) {0};
  \node[bit, fill=ln-tip!22] at (44.75mm,-24.0mm) {1};
  % FP8 E4M3
  \node[anchor=east] at (22mm,-30.0mm) {FP8 E4M3\enspace\textcolor{ln-muted}{1/4/3}};
  \node[bit, fill=ln-caution!30] at (24.45mm,-30.0mm) {0};
  \node[bit, fill=ln-accent!22] at (27.35mm,-30.0mm) {0};
  \node[bit, fill=ln-accent!22] at (30.25mm,-30.0mm) {1};
  \node[bit, fill=ln-accent!22] at (33.15mm,-30.0mm) {0};
  \node[bit, fill=ln-accent!22] at (36.05mm,-30.0mm) {0};
  \node[bit, fill=ln-tip!22] at (38.95mm,-30.0mm) {0};
  \node[bit, fill=ln-tip!22] at (41.85mm,-30.0mm) {1};
  \node[bit, fill=ln-tip!22] at (44.75mm,-30.0mm) {0};
  % legend
  \node[bit, fill=ln-caution!30] at (24.45mm,7mm) {};
  \node[anchor=west] at (26.45mm,7mm) {\iflangja{符号}{sign}};
  \node[bit, fill=ln-accent!22] at (44.45mm,7mm) {};
  \node[anchor=west] at (46.45mm,7mm) {\iflangja{指数部}{exponent}};
  \node[bit, fill=ln-tip!22] at (68.45mm,7mm) {};
  \node[anchor=west] at (70.45mm,7mm) {\iflangja{仮数部}{fraction}};
\end{tikzpicture}

% Width: body (120 mm); the drawing is about 116 mm wide.
\begin{tikzpicture}[font=\scriptsize\sffamily,
    bit/.style={draw=ln-muted, minimum width=2.9mm, minimum height=4.2mm,
      inner sep=0pt, font=\scriptsize\ttfamily}]
  % FP32
  \node[anchor=east] at (22mm,0.0mm) {FP32\enspace\textcolor{ln-muted}{1/8/23}};
  \node[bit, fill=ln-caution!30] at (24.45mm,0.0mm) {0};
  \node[bit, fill=ln-accent!22] at (27.35mm,0.0mm) {0};
  \node[bit, fill=ln-accent!22] at (30.25mm,0.0mm) {1};
  \node[bit, fill=ln-accent!22] at (33.15mm,0.0mm) {1};
  \node[bit, fill=ln-accent!22] at (36.05mm,0.0mm) {1};
  \node[bit, fill=ln-accent!22] at (38.95mm,0.0mm) {1};
  \node[bit, fill=ln-accent!22] at (41.85mm,0.0mm) {1};
  \node[bit, fill=ln-accent!22] at (44.75mm,0.0mm) {0};
  \node[bit, fill=ln-accent!22] at (47.65mm,0.0mm) {0};
  \node[bit, fill=ln-tip!22] at (50.55mm,0.0mm) {0};
  \node[bit, fill=ln-tip!22] at (53.45mm,0.0mm) {1};
  \node[bit, fill=ln-tip!22] at (56.35mm,0.0mm) {0};
  \node[bit, fill=ln-tip!22] at (59.25mm,0.0mm) {0};
  \node[bit, fill=ln-tip!22] at (62.15mm,0.0mm) {0};
  \node[bit, fill=ln-tip!22] at (65.05mm,0.0mm) {0};
  \node[bit, fill=ln-tip!22] at (67.95mm,0.0mm) {0};
  \node[bit, fill=ln-tip!22] at (70.85mm,0.0mm) {0};
  \node[bit, fill=ln-tip!22] at (73.75mm,0.0mm) {0};
  \node[bit, fill=ln-tip!22] at (76.65mm,0.0mm) {0};
  \node[bit, fill=ln-tip!22] at (79.55mm,0.0mm) {0};
  \node[bit, fill=ln-tip!22] at (82.45mm,0.0mm) {0};
  \node[bit, fill=ln-tip!22] at (85.35mm,0.0mm) {0};
  \node[bit, fill=ln-tip!22] at (88.25mm,0.0mm) {0};
  \node[bit, fill=ln-tip!22] at (91.15mm,0.0mm) {0};
  \node[bit, fill=ln-tip!22] at (94.05mm,0.0mm) {0};
  \node[bit, fill=ln-tip!22] at (96.95mm,0.0mm) {0};
  \node[bit, fill=ln-tip!22] at (99.85mm,0.0mm) {0};
  \node[bit, fill=ln-tip!22] at (102.75mm,0.0mm) {0};
  \node[bit, fill=ln-tip!22] at (105.65mm,0.0mm) {0};
  \node[bit, fill=ln-tip!22] at (108.55mm,0.0mm) {0};
  \node[bit, fill=ln-tip!22] at (111.45mm,0.0mm) {0};
  \node[bit, fill=ln-tip!22] at (114.35mm,0.0mm) {0};
  % TF32
  \node[anchor=east] at (22mm,-6.0mm) {TF32\enspace\textcolor{ln-muted}{1/8/10}};
  \node[bit, fill=ln-caution!30] at (24.45mm,-6.0mm) {0};
  \node[bit, fill=ln-accent!22] at (27.35mm,-6.0mm) {0};
  \node[bit, fill=ln-accent!22] at (30.25mm,-6.0mm) {1};
  \node[bit, fill=ln-accent!22] at (33.15mm,-6.0mm) {1};
  \node[bit, fill=ln-accent!22] at (36.05mm,-6.0mm) {1};
  \node[bit, fill=ln-accent!22] at (38.95mm,-6.0mm) {1};
  \node[bit, fill=ln-accent!22] at (41.85mm,-6.0mm) {1};
  \node[bit, fill=ln-accent!22] at (44.75mm,-6.0mm) {0};
  \node[bit, fill=ln-accent!22] at (47.65mm,-6.0mm) {0};
  \node[bit, fill=ln-tip!22] at (50.55mm,-6.0mm) {0};
  \node[bit, fill=ln-tip!22] at (53.45mm,-6.0mm) {1};
  \node[bit, fill=ln-tip!22] at (56.35mm,-6.0mm) {0};
  \node[bit, fill=ln-tip!22] at (59.25mm,-6.0mm) {0};
  \node[bit, fill=ln-tip!22] at (62.15mm,-6.0mm) {0};
  \node[bit, fill=ln-tip!22] at (65.05mm,-6.0mm) {0};
  \node[bit, fill=ln-tip!22] at (67.95mm,-6.0mm) {0};
  \node[bit, fill=ln-tip!22] at (70.85mm,-6.0mm) {0};
  \node[bit, fill=ln-tip!22] at (73.75mm,-6.0mm) {0};
  \node[bit, fill=ln-tip!22] at (76.65mm,-6.0mm) {0};
  % BF16
  \node[anchor=east] at (22mm,-12.0mm) {BF16\enspace\textcolor{ln-muted}{1/8/7}};
  \node[bit, fill=ln-caution!30] at (24.45mm,-12.0mm) {0};
  \node[bit, fill=ln-accent!22] at (27.35mm,-12.0mm) {0};
  \node[bit, fill=ln-accent!22] at (30.25mm,-12.0mm) {1};
  \node[bit, fill=ln-accent!22] at (33.15mm,-12.0mm) {1};
  \node[bit, fill=ln-accent!22] at (36.05mm,-12.0mm) {1};
  \node[bit, fill=ln-accent!22] at (38.95mm,-12.0mm) {1};
  \node[bit, fill=ln-accent!22] at (41.85mm,-12.0mm) {1};
  \node[bit, fill=ln-accent!22] at (44.75mm,-12.0mm) {0};
  \node[bit, fill=ln-accent!22] at (47.65mm,-12.0mm) {0};
  \node[bit, fill=ln-tip!22] at (50.55mm,-12.0mm) {0};
  \node[bit, fill=ln-tip!22] at (53.45mm,-12.0mm) {1};
  \node[bit, fill=ln-tip!22] at (56.35mm,-12.0mm) {0};
  \node[bit, fill=ln-tip!22] at (59.25mm,-12.0mm) {0};
  \node[bit, fill=ln-tip!22] at (62.15mm,-12.0mm) {0};
  \node[bit, fill=ln-tip!22] at (65.05mm,-12.0mm) {0};
  \node[bit, fill=ln-tip!22] at (67.95mm,-12.0mm) {0};
  % FP16
  \node[anchor=east] at (22mm,-18.0mm) {FP16\enspace\textcolor{ln-muted}{1/5/10}};
  \node[bit, fill=ln-caution!30] at (24.45mm,-18.0mm) {0};
  \node[bit, fill=ln-accent!22] at (27.35mm,-18.0mm) {0};
  \node[bit, fill=ln-accent!22] at (30.25mm,-18.0mm) {1};
  \node[bit, fill=ln-accent!22] at (33.15mm,-18.0mm) {1};
  \node[bit, fill=ln-accent!22] at (36.05mm,-18.0mm) {0};
  \node[bit, fill=ln-accent!22] at (38.95mm,-18.0mm) {0};
  \node[bit, fill=ln-tip!22] at (41.85mm,-18.0mm) {0};
  \node[bit, fill=ln-tip!22] at (44.75mm,-18.0mm) {1};
  \node[bit, fill=ln-tip!22] at (47.65mm,-18.0mm) {0};
  \node[bit, fill=ln-tip!22] at (50.55mm,-18.0mm) {0};
  \node[bit, fill=ln-tip!22] at (53.45mm,-18.0mm) {0};
  \node[bit, fill=ln-tip!22] at (56.35mm,-18.0mm) {0};
  \node[bit, fill=ln-tip!22] at (59.25mm,-18.0mm) {0};
  \node[bit, fill=ln-tip!22] at (62.15mm,-18.0mm) {0};
  \node[bit, fill=ln-tip!22] at (65.05mm,-18.0mm) {0};
  \node[bit, fill=ln-tip!22] at (67.95mm,-18.0mm) {0};
  % FP8 E5M2
  \node[anchor=east] at (22mm,-24.0mm) {FP8 E5M2\enspace\textcolor{ln-muted}{1/5/2}};
  \node[bit, fill=ln-caution!30] at (24.45mm,-24.0mm) {0};
  \node[bit, fill=ln-accent!22] at (27.35mm,-24.0mm) {0};
  \node[bit, fill=ln-accent!22] at (30.25mm,-24.0mm) {1};
  \node[bit, fill=ln-accent!22] at (33.15mm,-24.0mm) {1};
  \node[bit, fill=ln-accent!22] at (36.05mm,-24.0mm) {0};
  \node[bit, fill=ln-accent!22] at (38.95mm,-24.0mm) {0};
  \node[bit, fill=ln-tip!22] at (41.85mm,-24.0mm) {0};
  \node[bit, fill=ln-tip!22] at (44.75mm,-24.0mm) {1};
  % FP8 E4M3
  \node[anchor=east] at (22mm,-30.0mm) {FP8 E4M3\enspace\textcolor{ln-muted}{1/4/3}};
  \node[bit, fill=ln-caution!30] at (24.45mm,-30.0mm) {0};
  \node[bit, fill=ln-accent!22] at (27.35mm,-30.0mm) {0};
  \node[bit, fill=ln-accent!22] at (30.25mm,-30.0mm) {1};
  \node[bit, fill=ln-accent!22] at (33.15mm,-30.0mm) {0};
  \node[bit, fill=ln-accent!22] at (36.05mm,-30.0mm) {0};
  \node[bit, fill=ln-tip!22] at (38.95mm,-30.0mm) {0};
  \node[bit, fill=ln-tip!22] at (41.85mm,-30.0mm) {1};
  \node[bit, fill=ln-tip!22] at (44.75mm,-30.0mm) {0};
  % legend
  \node[bit, fill=ln-caution!30] at (24.45mm,7mm) {};
  \node[anchor=west] at (26.45mm,7mm) {\iflangja{符号}{sign}};
  \node[bit, fill=ln-accent!22] at (44.45mm,7mm) {};
  \node[anchor=west] at (46.45mm,7mm) {\iflangja{指数部}{exponent}};
  \node[bit, fill=ln-tip!22] at (68.45mm,7mm) {};
  \node[anchor=west] at (70.45mm,7mm) {\iflangja{仮数部}{fraction}};
\end{tikzpicture}

% Width: body (120 mm); the drawing is about 116 mm wide.
\begin{tikzpicture}[font=\scriptsize\sffamily,
    bit/.style={draw=ln-muted, minimum width=2.9mm, minimum height=4.2mm,
      inner sep=0pt, font=\scriptsize\ttfamily}]
  % FP32
  \node[anchor=east] at (22mm,0.0mm) {FP32\enspace\textcolor{ln-muted}{1/8/23}};
  \node[bit, fill=ln-caution!30] at (24.45mm,0.0mm) {0};
  \node[bit, fill=ln-accent!22] at (27.35mm,0.0mm) {0};
  \node[bit, fill=ln-accent!22] at (30.25mm,0.0mm) {1};
  \node[bit, fill=ln-accent!22] at (33.15mm,0.0mm) {1};
  \node[bit, fill=ln-accent!22] at (36.05mm,0.0mm) {1};
  \node[bit, fill=ln-accent!22] at (38.95mm,0.0mm) {1};
  \node[bit, fill=ln-accent!22] at (41.85mm,0.0mm) {1};
  \node[bit, fill=ln-accent!22] at (44.75mm,0.0mm) {0};
  \node[bit, fill=ln-accent!22] at (47.65mm,0.0mm) {0};
  \node[bit, fill=ln-tip!22] at (50.55mm,0.0mm) {0};
  \node[bit, fill=ln-tip!22] at (53.45mm,0.0mm) {1};
  \node[bit, fill=ln-tip!22] at (56.35mm,0.0mm) {0};
  \node[bit, fill=ln-tip!22] at (59.25mm,0.0mm) {0};
  \node[bit, fill=ln-tip!22] at (62.15mm,0.0mm) {0};
  \node[bit, fill=ln-tip!22] at (65.05mm,0.0mm) {0};
  \node[bit, fill=ln-tip!22] at (67.95mm,0.0mm) {0};
  \node[bit, fill=ln-tip!22] at (70.85mm,0.0mm) {0};
  \node[bit, fill=ln-tip!22] at (73.75mm,0.0mm) {0};
  \node[bit, fill=ln-tip!22] at (76.65mm,0.0mm) {0};
  \node[bit, fill=ln-tip!22] at (79.55mm,0.0mm) {0};
  \node[bit, fill=ln-tip!22] at (82.45mm,0.0mm) {0};
  \node[bit, fill=ln-tip!22] at (85.35mm,0.0mm) {0};
  \node[bit, fill=ln-tip!22] at (88.25mm,0.0mm) {0};
  \node[bit, fill=ln-tip!22] at (91.15mm,0.0mm) {0};
  \node[bit, fill=ln-tip!22] at (94.05mm,0.0mm) {0};
  \node[bit, fill=ln-tip!22] at (96.95mm,0.0mm) {0};
  \node[bit, fill=ln-tip!22] at (99.85mm,0.0mm) {0};
  \node[bit, fill=ln-tip!22] at (102.75mm,0.0mm) {0};
  \node[bit, fill=ln-tip!22] at (105.65mm,0.0mm) {0};
  \node[bit, fill=ln-tip!22] at (108.55mm,0.0mm) {0};
  \node[bit, fill=ln-tip!22] at (111.45mm,0.0mm) {0};
  \node[bit, fill=ln-tip!22] at (114.35mm,0.0mm) {0};
  % TF32
  \node[anchor=east] at (22mm,-6.0mm) {TF32\enspace\textcolor{ln-muted}{1/8/10}};
  \node[bit, fill=ln-caution!30] at (24.45mm,-6.0mm) {0};
  \node[bit, fill=ln-accent!22] at (27.35mm,-6.0mm) {0};
  \node[bit, fill=ln-accent!22] at (30.25mm,-6.0mm) {1};
  \node[bit, fill=ln-accent!22] at (33.15mm,-6.0mm) {1};
  \node[bit, fill=ln-accent!22] at (36.05mm,-6.0mm) {1};
  \node[bit, fill=ln-accent!22] at (38.95mm,-6.0mm) {1};
  \node[bit, fill=ln-accent!22] at (41.85mm,-6.0mm) {1};
  \node[bit, fill=ln-accent!22] at (44.75mm,-6.0mm) {0};
  \node[bit, fill=ln-accent!22] at (47.65mm,-6.0mm) {0};
  \node[bit, fill=ln-tip!22] at (50.55mm,-6.0mm) {0};
  \node[bit, fill=ln-tip!22] at (53.45mm,-6.0mm) {1};
  \node[bit, fill=ln-tip!22] at (56.35mm,-6.0mm) {0};
  \node[bit, fill=ln-tip!22] at (59.25mm,-6.0mm) {0};
  \node[bit, fill=ln-tip!22] at (62.15mm,-6.0mm) {0};
  \node[bit, fill=ln-tip!22] at (65.05mm,-6.0mm) {0};
  \node[bit, fill=ln-tip!22] at (67.95mm,-6.0mm) {0};
  \node[bit, fill=ln-tip!22] at (70.85mm,-6.0mm) {0};
  \node[bit, fill=ln-tip!22] at (73.75mm,-6.0mm) {0};
  \node[bit, fill=ln-tip!22] at (76.65mm,-6.0mm) {0};
  % BF16
  \node[anchor=east] at (22mm,-12.0mm) {BF16\enspace\textcolor{ln-muted}{1/8/7}};
  \node[bit, fill=ln-caution!30] at (24.45mm,-12.0mm) {0};
  \node[bit, fill=ln-accent!22] at (27.35mm,-12.0mm) {0};
  \node[bit, fill=ln-accent!22] at (30.25mm,-12.0mm) {1};
  \node[bit, fill=ln-accent!22] at (33.15mm,-12.0mm) {1};
  \node[bit, fill=ln-accent!22] at (36.05mm,-12.0mm) {1};
  \node[bit, fill=ln-accent!22] at (38.95mm,-12.0mm) {1};
  \node[bit, fill=ln-accent!22] at (41.85mm,-12.0mm) {1};
  \node[bit, fill=ln-accent!22] at (44.75mm,-12.0mm) {0};
  \node[bit, fill=ln-accent!22] at (47.65mm,-12.0mm) {0};
  \node[bit, fill=ln-tip!22] at (50.55mm,-12.0mm) {0};
  \node[bit, fill=ln-tip!22] at (53.45mm,-12.0mm) {1};
  \node[bit, fill=ln-tip!22] at (56.35mm,-12.0mm) {0};
  \node[bit, fill=ln-tip!22] at (59.25mm,-12.0mm) {0};
  \node[bit, fill=ln-tip!22] at (62.15mm,-12.0mm) {0};
  \node[bit, fill=ln-tip!22] at (65.05mm,-12.0mm) {0};
  \node[bit, fill=ln-tip!22] at (67.95mm,-12.0mm) {0};
  % FP16
  \node[anchor=east] at (22mm,-18.0mm) {FP16\enspace\textcolor{ln-muted}{1/5/10}};
  \node[bit, fill=ln-caution!30] at (24.45mm,-18.0mm) {0};
  \node[bit, fill=ln-accent!22] at (27.35mm,-18.0mm) {0};
  \node[bit, fill=ln-accent!22] at (30.25mm,-18.0mm) {1};
  \node[bit, fill=ln-accent!22] at (33.15mm,-18.0mm) {1};
  \node[bit, fill=ln-accent!22] at (36.05mm,-18.0mm) {0};
  \node[bit, fill=ln-accent!22] at (38.95mm,-18.0mm) {0};
  \node[bit, fill=ln-tip!22] at (41.85mm,-18.0mm) {0};
  \node[bit, fill=ln-tip!22] at (44.75mm,-18.0mm) {1};
  \node[bit, fill=ln-tip!22] at (47.65mm,-18.0mm) {0};
  \node[bit, fill=ln-tip!22] at (50.55mm,-18.0mm) {0};
  \node[bit, fill=ln-tip!22] at (53.45mm,-18.0mm) {0};
  \node[bit, fill=ln-tip!22] at (56.35mm,-18.0mm) {0};
  \node[bit, fill=ln-tip!22] at (59.25mm,-18.0mm) {0};
  \node[bit, fill=ln-tip!22] at (62.15mm,-18.0mm) {0};
  \node[bit, fill=ln-tip!22] at (65.05mm,-18.0mm) {0};
  \node[bit, fill=ln-tip!22] at (67.95mm,-18.0mm) {0};
  % FP8 E5M2
  \node[anchor=east] at (22mm,-24.0mm) {FP8 E5M2\enspace\textcolor{ln-muted}{1/5/2}};
  \node[bit, fill=ln-caution!30] at (24.45mm,-24.0mm) {0};
  \node[bit, fill=ln-accent!22] at (27.35mm,-24.0mm) {0};
  \node[bit, fill=ln-accent!22] at (30.25mm,-24.0mm) {1};
  \node[bit, fill=ln-accent!22] at (33.15mm,-24.0mm) {1};
  \node[bit, fill=ln-accent!22] at (36.05mm,-24.0mm) {0};
  \node[bit, fill=ln-accent!22] at (38.95mm,-24.0mm) {0};
  \node[bit, fill=ln-tip!22] at (41.85mm,-24.0mm) {0};
  \node[bit, fill=ln-tip!22] at (44.75mm,-24.0mm) {1};
  % FP8 E4M3
  \node[anchor=east] at (22mm,-30.0mm) {FP8 E4M3\enspace\textcolor{ln-muted}{1/4/3}};
  \node[bit, fill=ln-caution!30] at (24.45mm,-30.0mm) {0};
  \node[bit, fill=ln-accent!22] at (27.35mm,-30.0mm) {0};
  \node[bit, fill=ln-accent!22] at (30.25mm,-30.0mm) {1};
  \node[bit, fill=ln-accent!22] at (33.15mm,-30.0mm) {0};
  \node[bit, fill=ln-accent!22] at (36.05mm,-30.0mm) {0};
  \node[bit, fill=ln-tip!22] at (38.95mm,-30.0mm) {0};
  \node[bit, fill=ln-tip!22] at (41.85mm,-30.0mm) {1};
  \node[bit, fill=ln-tip!22] at (44.75mm,-30.0mm) {0};
  % legend
  \node[bit, fill=ln-caution!30] at (24.45mm,7mm) {};
  \node[anchor=west] at (26.45mm,7mm) {\iflangja{符号}{sign}};
  \node[bit, fill=ln-accent!22] at (44.45mm,7mm) {};
  \node[anchor=west] at (46.45mm,7mm) {\iflangja{指数部}{exponent}};
  \node[bit, fill=ln-tip!22] at (68.45mm,7mm) {};
  \node[anchor=west] at (70.45mm,7mm) {\iflangja{仮数部}{fraction}};
\end{tikzpicture}

  \caption{The value $0.15625 = 1.01_2 \times 2^{-3}$ in six formats
    (sign/exponent/fraction bits on the left).  The biased exponent is
    $-3 + 127 = 124$ in FP32, TF32 and BF16, $-3 + 15 = 12$ in FP16 and
    E5M2, and $-3 + 7 = 4$ in E4M3; the fraction is $01$ followed by
    zeros.}
  \label{fig:l12-formats}
\end{figure}

\begin{definition}[ML formats]
  The format \term{bfloat16 (BF16)}[bfloat16] \enquote{includes an 8-bit exponent,
  7-bit mantissa, and 1 sign-bit} \cite{nvidia2020ampere}: the range of
  FP32 with less precision than FP16.\margin{In practice E4M3 tends to be used for
weights and activations and E5M2 for gradients, which span a wider
range \cite{micikevicius2022fp8}.}
  The format \term{TensorFloat-32 (TF32)}[TF32] has an \enquote{8-bit exponent
  (same as FP32), 10-bit mantissa (same precision as FP16) and 1
  sign-bit} \cite{nvidia2020ampere}; the A100 tensor cores use it for
  FP32 inputs.  The 8-bit format \term{FP8}[FP8] has two encodings: \enquote{E4M3 with 4
  exponent bits, 3 mantissa bits, and 1 sign bit} and \enquote{E5M2,
  with 5 exponent bits, 2 mantissa bits, and 1 sign bit}
  \cite{nvidia2022hopper}.
\end{definition}

The two FP8 encodings divide the work of training: \enquote{E4M3
supports computations requiring less dynamic range with more precision,
while E5M2 provides a wider dynamic range and less precision}
\cite{nvidia2022hopper}.\caution{E4M3 is not IEEE-like at the top: it has no infinity, and only
the all-ones pattern \texttt{S.1111.111} is NaN \cite{micikevicius2022fp8};
the PTX ISA lists its NaNs as \texttt{0x7f} and \texttt{0xff}
\cite[§5.2.3 Alternate Floating-Point Data Formats]{nvidia2026ptx}.  The
largest value is therefore $1.110_2 \times 2^{15-7} = 1.75 \times 256 =
448$, not the $240$ the IEEE rules would give.}  The
\term{dynamic range}[ダイナミックレンジ] of a format is the ratio of
its largest to its smallest representable positive magnitude, often
quoted in powers of two (binades).

\begin{table}[htbp]
  \centering
  \caption{Parameters of the formats of \cref{fig:l12-formats}.  Maxima
    and minima are derived from \cref{eq:l12-normal}, except the E4M3
    maximum, which uses the exponent field $1111$ (no infinities; see
    the caution in the text); $\varepsilon =
    2^{-p}$ is the spacing of values just above~1.}
  \label{tab:l12-formats}
  \setlength{\tabcolsep}{4pt}
  \begin{tabular}{lcrlll}
    \toprule
    Format & $s/w/p$ & Bias & Max normal & Min normal & $\varepsilon$ \\
    \midrule
    FP64 & 1/11/52 & 1023 & $\approx 1.80 \times 10^{308}$ & $2^{-1022}$ & $2^{-52}$ \\
    FP32 & 1/8/23 & 127 & $\approx 3.40 \times 10^{38}$ & $2^{-126}$ & $2^{-23}$ \\
    TF32 & 1/8/10 & 127 & $\approx 3.40 \times 10^{38}$ & $2^{-126}$ & $2^{-10}$ \\
    BF16 & 1/8/7 & 127 & $\approx 3.39 \times 10^{38}$ & $2^{-126}$ & $2^{-7}$ \\
    \addlinespace
    FP16 & 1/5/10 & 15 & \num{65504} & $2^{-14}$ & $2^{-10}$ \\
    FP8 E5M2 & 1/5/2 & 15 & \num{57344} & $2^{-14}$ & $2^{-2}$ \\
    FP8 E4M3 & 1/4/3 & 7 & 448 & $2^{-6}$ & $2^{-3}$ \\
    \bottomrule
  \end{tabular}
\end{table}



\begin{example}
  \label{ex:l12-bf16}
  Range against precision in FP16 and BF16.\tip{BF16 is the upper
  half of FP32: converting FP32 to BF16 drops the low 16 fraction bits,
  by truncation or by rounding.}  The value \num{70000}
  exceeds the FP16 maximum of \num{65504} and overflows; in BF16 it is
  in range, and since the spacing of BF16 values between $2^{16}$ and
  $2^{17}$ is $2^{16-7} = 512$, it rounds to the nearest multiple,
  $137 \times 512 = \num{70144}$.  Conversely, $1 + 2^{-8}$ is exact in
  FP16 ($\varepsilon = 2^{-10}$), but in BF16 it lies halfway between
  $1$ and $1 + 2^{-7}$ and rounds to the neighbour with an even
  fraction, $1$: the increment is lost.
\end{example}

\subsection{Subnormals, special values and IEEE 754}

\Cref{eq:l12-normal} cannot represent zero, and its smallest magnitude
$2^{1-b}$ leaves a gap down to zero.  The exponent field values $0$ and
$2^w - 1$ are reserved to fill these holes.

\begin{definition}[Subnormal number, underflow, overflow, NaN]
  A \term{subnormal number}[非正規化数] has exponent field $e = 0$ and
  the value $(-1)^{s} \times (f / 2^{p}) \times 2^{1 - b}$, without the
  implicit leading~1; with $f = 0$ it is $\pm 0$.  Subnormals fill the
  gap between zero and the smallest normal number with evenly spaced
  values.  A nonzero result whose magnitude is below the smallest normal number
  $2^{1-b}$ suffers \term{underflow}[アンダーフロー] (it is rounded to a
  subnormal or to zero); a result beyond the largest finite value suffers
  \term{overflow}[オーバーフロー] (under the default rounding it becomes
  $\pm\infty$; a format without infinities, such as E4M3, saturates or
  gives NaN).  The exponent
  field $e = 2^w - 1$ encodes $\pm\infty$ ($f = 0$) and
  \term{NaN (not a number)}[非数（NaN）] ($f \ne 0$), the result of an
  invalid operation such as $0/0$ or $\infty - \infty$.
\end{definition}

\Cref{fig:l12-number-line}(a) lists every non-negative value of a toy
format with two exponent bits (bias~2) and two fraction bits,\margin{Toy format: $s/w/p = 1/2/2$, $b = 2$.
Normals: $1.f \times 2^{e-2}$ for $e = 1, 2, 3$, from $0.5$ to $3.5$;
subnormals $0.f \times 2^{-1}$: $1/8$, $2/8$, $3/8$.} using all
exponent codes for numbers.  The normal values in $[2^{k}, 2^{k+1})$ are
evenly spaced, and the spacing doubles at each power of two; the three
subnormals continue the spacing of the lowest binade down to zero.
Representable values are therefore dense near zero and sparse at large
magnitudes: precision, the spacing relative to the value, is roughly
constant, but the absolute error grows with the magnitude.
\Cref{fig:l12-number-line}(b) shows the regions of a real format on a
logarithmic scale.

\begin{figure}[tb]
  \centering
  % (a) All non-negative values of a toy floating-point format with 1 sign
% bit, 2 exponent bits (bias 2) and 2 fraction bits, without infinities:
% subnormals 1/8, 2/8, 3/8 (exponent field 00), normals 0.5 ... 3.5, whose
% spacing doubles at every power of two.  (b) The regions of a real format
% on a logarithmic magnitude axis: overflow, normal, subnormal and
% underflow on each side of zero.
% Redraws the outline figures m12l2_floating_basics_diff_details (number
% line) and m12l2_floating_rep (regions).
% Usage: % (a) All non-negative values of a toy floating-point format with 1 sign
% bit, 2 exponent bits (bias 2) and 2 fraction bits, without infinities:
% subnormals 1/8, 2/8, 3/8 (exponent field 00), normals 0.5 ... 3.5, whose
% spacing doubles at every power of two.  (b) The regions of a real format
% on a logarithmic magnitude axis: overflow, normal, subnormal and
% underflow on each side of zero.
% Redraws the outline figures m12l2_floating_basics_diff_details (number
% line) and m12l2_floating_rep (regions).
% Usage: % (a) All non-negative values of a toy floating-point format with 1 sign
% bit, 2 exponent bits (bias 2) and 2 fraction bits, without infinities:
% subnormals 1/8, 2/8, 3/8 (exponent field 00), normals 0.5 ... 3.5, whose
% spacing doubles at every power of two.  (b) The regions of a real format
% on a logarithmic magnitude axis: overflow, normal, subnormal and
% underflow on each side of zero.
% Redraws the outline figures m12l2_floating_basics_diff_details (number
% line) and m12l2_floating_rep (regions).
% Usage: \input{figures/l12-number-line}
% Width: body (120 mm); the drawing is about 112 mm wide.
\begin{tikzpicture}[font=\scriptsize\sffamily, x=26mm,
    >={Stealth[length=1.3mm]}]
  % (a) number line 0..4
  \node[anchor=west] at (-0.3,1.05) {(a)};
  \draw[->] (0,0) -- (4.15,0);
  \foreach \v in {0,1,2,3,4} {
    \draw (\v,0.12) -- (\v,-0.12) node[below=1pt] {\v};
  }
  \fill (0,0) circle (0.5mm);
  \foreach \v in {0.125,0.25,0.375} {
    \fill[ln-caution] (\v,0) circle (0.55mm);
  }
  \foreach \v in {0.5,0.625,0.75,0.875,1,1.25,1.5,1.75,2,2.5,3,3.5} {
    \fill[ln-accent] (\v,0) circle (0.55mm);
  }
  \node[below=1pt] at (0.5,-0.12) {0.5};
  \node[below=1pt] at (3.5,-0.12) {3.5};
  % spacing annotations
  \draw[|-|, ln-muted] (0,0.35) -- node[above] {1/8} (0.125,0.35);
  \draw[|-|, ln-muted] (0.5,0.35) -- node[above] {1/8} (0.625,0.35);
  \draw[|-|, ln-muted] (1,0.35) -- node[above] {1/4} (1.25,0.35);
  \draw[|-|, ln-muted] (2,0.35) -- node[above] {1/2} (2.5,0.35);
  \node[text=ln-caution, anchor=west] at (0.0,-0.75)
    {\iflangja{● 非正規化数（指数フィールド 00）}{$\bullet$ subnormal (exponent field 00)}};
  \node[text=ln-accent, anchor=west] at (2.0,-0.75)
    {\iflangja{● 正規化数（01, 10, 11）}{$\bullet$ normal (fields 01, 10, 11)}};
  % (b) regions, log scale
  \begin{scope}[yshift=-24mm]
    \node[anchor=west] at (-0.3,0.75) {(b)};
    \def\h{0.42}
    \foreach \a/\b/\c/\t in {
        -0.2/0.3/ln-caution!25/{\iflangja{オーバー\\フロー}{over-\\flow}},
        0.3/1.25/ln-accent!18/{\iflangja{正規化数}{normal}},
        1.25/1.6/ln-tip!25/{\iflangja{非正規}{sub-\\normal}},
        1.6/1.85/ln-muted!20/{},
        2.15/2.4/ln-muted!20/{},
        2.4/2.75/ln-tip!25/{\iflangja{非正規}{sub-\\normal}},
        2.75/3.7/ln-accent!18/{\iflangja{正規化数}{normal}},
        3.7/4.2/ln-caution!25/{\iflangja{オーバー\\フロー}{over-\\flow}}} {
      \draw[fill=\c] (\a,-\h) rectangle (\b,\h);
      \node[align=center, inner sep=0pt] at ({(\a+\b)/2},0) {\t};
    }
    \node at (2,0) {0};
    \draw[ln-muted] (1.85,-\h) rectangle (2.15,\h);
    \draw[->, ln-muted] (1.72,-0.6) -- (1.72,-\h);
    \draw[->, ln-muted] (2.28,-0.6) -- (2.28,-\h);
    \node[anchor=north, text=ln-muted] at (2,-0.6)
      {\iflangja{アンダーフロー}{underflow}};
    \draw[->, ln-muted] (1.6,0.65) -- node[above]
      {\iflangja{負，絶対値が大}{negative, larger magnitude}} (0.0,0.65);
    \draw[->, ln-muted] (2.4,0.65) -- node[above]
      {\iflangja{正，絶対値が大}{positive, larger magnitude}} (4.0,0.65);
  \end{scope}
\end{tikzpicture}

% Width: body (120 mm); the drawing is about 112 mm wide.
\begin{tikzpicture}[font=\scriptsize\sffamily, x=26mm,
    >={Stealth[length=1.3mm]}]
  % (a) number line 0..4
  \node[anchor=west] at (-0.3,1.05) {(a)};
  \draw[->] (0,0) -- (4.15,0);
  \foreach \v in {0,1,2,3,4} {
    \draw (\v,0.12) -- (\v,-0.12) node[below=1pt] {\v};
  }
  \fill (0,0) circle (0.5mm);
  \foreach \v in {0.125,0.25,0.375} {
    \fill[ln-caution] (\v,0) circle (0.55mm);
  }
  \foreach \v in {0.5,0.625,0.75,0.875,1,1.25,1.5,1.75,2,2.5,3,3.5} {
    \fill[ln-accent] (\v,0) circle (0.55mm);
  }
  \node[below=1pt] at (0.5,-0.12) {0.5};
  \node[below=1pt] at (3.5,-0.12) {3.5};
  % spacing annotations
  \draw[|-|, ln-muted] (0,0.35) -- node[above] {1/8} (0.125,0.35);
  \draw[|-|, ln-muted] (0.5,0.35) -- node[above] {1/8} (0.625,0.35);
  \draw[|-|, ln-muted] (1,0.35) -- node[above] {1/4} (1.25,0.35);
  \draw[|-|, ln-muted] (2,0.35) -- node[above] {1/2} (2.5,0.35);
  \node[text=ln-caution, anchor=west] at (0.0,-0.75)
    {\iflangja{● 非正規化数（指数フィールド 00）}{$\bullet$ subnormal (exponent field 00)}};
  \node[text=ln-accent, anchor=west] at (2.0,-0.75)
    {\iflangja{● 正規化数（01, 10, 11）}{$\bullet$ normal (fields 01, 10, 11)}};
  % (b) regions, log scale
  \begin{scope}[yshift=-24mm]
    \node[anchor=west] at (-0.3,0.75) {(b)};
    \def\h{0.42}
    \foreach \a/\b/\c/\t in {
        -0.2/0.3/ln-caution!25/{\iflangja{オーバー\\フロー}{over-\\flow}},
        0.3/1.25/ln-accent!18/{\iflangja{正規化数}{normal}},
        1.25/1.6/ln-tip!25/{\iflangja{非正規}{sub-\\normal}},
        1.6/1.85/ln-muted!20/{},
        2.15/2.4/ln-muted!20/{},
        2.4/2.75/ln-tip!25/{\iflangja{非正規}{sub-\\normal}},
        2.75/3.7/ln-accent!18/{\iflangja{正規化数}{normal}},
        3.7/4.2/ln-caution!25/{\iflangja{オーバー\\フロー}{over-\\flow}}} {
      \draw[fill=\c] (\a,-\h) rectangle (\b,\h);
      \node[align=center, inner sep=0pt] at ({(\a+\b)/2},0) {\t};
    }
    \node at (2,0) {0};
    \draw[ln-muted] (1.85,-\h) rectangle (2.15,\h);
    \draw[->, ln-muted] (1.72,-0.6) -- (1.72,-\h);
    \draw[->, ln-muted] (2.28,-0.6) -- (2.28,-\h);
    \node[anchor=north, text=ln-muted] at (2,-0.6)
      {\iflangja{アンダーフロー}{underflow}};
    \draw[->, ln-muted] (1.6,0.65) -- node[above]
      {\iflangja{負，絶対値が大}{negative, larger magnitude}} (0.0,0.65);
    \draw[->, ln-muted] (2.4,0.65) -- node[above]
      {\iflangja{正，絶対値が大}{positive, larger magnitude}} (4.0,0.65);
  \end{scope}
\end{tikzpicture}

% Width: body (120 mm); the drawing is about 112 mm wide.
\begin{tikzpicture}[font=\scriptsize\sffamily, x=26mm,
    >={Stealth[length=1.3mm]}]
  % (a) number line 0..4
  \node[anchor=west] at (-0.3,1.05) {(a)};
  \draw[->] (0,0) -- (4.15,0);
  \foreach \v in {0,1,2,3,4} {
    \draw (\v,0.12) -- (\v,-0.12) node[below=1pt] {\v};
  }
  \fill (0,0) circle (0.5mm);
  \foreach \v in {0.125,0.25,0.375} {
    \fill[ln-caution] (\v,0) circle (0.55mm);
  }
  \foreach \v in {0.5,0.625,0.75,0.875,1,1.25,1.5,1.75,2,2.5,3,3.5} {
    \fill[ln-accent] (\v,0) circle (0.55mm);
  }
  \node[below=1pt] at (0.5,-0.12) {0.5};
  \node[below=1pt] at (3.5,-0.12) {3.5};
  % spacing annotations
  \draw[|-|, ln-muted] (0,0.35) -- node[above] {1/8} (0.125,0.35);
  \draw[|-|, ln-muted] (0.5,0.35) -- node[above] {1/8} (0.625,0.35);
  \draw[|-|, ln-muted] (1,0.35) -- node[above] {1/4} (1.25,0.35);
  \draw[|-|, ln-muted] (2,0.35) -- node[above] {1/2} (2.5,0.35);
  \node[text=ln-caution, anchor=west] at (0.0,-0.75)
    {\iflangja{● 非正規化数（指数フィールド 00）}{$\bullet$ subnormal (exponent field 00)}};
  \node[text=ln-accent, anchor=west] at (2.0,-0.75)
    {\iflangja{● 正規化数（01, 10, 11）}{$\bullet$ normal (fields 01, 10, 11)}};
  % (b) regions, log scale
  \begin{scope}[yshift=-24mm]
    \node[anchor=west] at (-0.3,0.75) {(b)};
    \def\h{0.42}
    \foreach \a/\b/\c/\t in {
        -0.2/0.3/ln-caution!25/{\iflangja{オーバー\\フロー}{over-\\flow}},
        0.3/1.25/ln-accent!18/{\iflangja{正規化数}{normal}},
        1.25/1.6/ln-tip!25/{\iflangja{非正規}{sub-\\normal}},
        1.6/1.85/ln-muted!20/{},
        2.15/2.4/ln-muted!20/{},
        2.4/2.75/ln-tip!25/{\iflangja{非正規}{sub-\\normal}},
        2.75/3.7/ln-accent!18/{\iflangja{正規化数}{normal}},
        3.7/4.2/ln-caution!25/{\iflangja{オーバー\\フロー}{over-\\flow}}} {
      \draw[fill=\c] (\a,-\h) rectangle (\b,\h);
      \node[align=center, inner sep=0pt] at ({(\a+\b)/2},0) {\t};
    }
    \node at (2,0) {0};
    \draw[ln-muted] (1.85,-\h) rectangle (2.15,\h);
    \draw[->, ln-muted] (1.72,-0.6) -- (1.72,-\h);
    \draw[->, ln-muted] (2.28,-0.6) -- (2.28,-\h);
    \node[anchor=north, text=ln-muted] at (2,-0.6)
      {\iflangja{アンダーフロー}{underflow}};
    \draw[->, ln-muted] (1.6,0.65) -- node[above]
      {\iflangja{負，絶対値が大}{negative, larger magnitude}} (0.0,0.65);
    \draw[->, ln-muted] (2.4,0.65) -- node[above]
      {\iflangja{正，絶対値が大}{positive, larger magnitude}} (4.0,0.65);
  \end{scope}
\end{tikzpicture}

  \caption{(a) The non-negative values of a toy format (1/2/2 bits, bias
    2, no infinities): three subnormals, then normals whose spacing
    doubles at every power of two.  (b) The regions of a format on a
    logarithmic magnitude axis.}
  \label{fig:l12-number-line}
\end{figure}

\begin{definition}[IEEE 754; rounding mode]
  The standard \term{IEEE 754}[IEEE 754] defines binary floating-point
  arithmetic \cite{ieee754-2019}.\margin{IEEE 754 binary16, binary32 and binary64 are the FP16,
FP32 and FP64 of this chapter (Lesson~1 introduced the three
precisions).}  It specifies the layout of sign,
  exponent and significand for its binary formats (binary16, binary32,
  binary64, \ldots), the special values ($\pm 0$, $\pm\infty$, NaN),
  subnormal numbers, the exact results of the basic operations, and the
  \term{rounding mode}[丸めモード]s that choose a representable result
  when the exact one is not representable: round toward $+\infty$
  (up), toward $-\infty$ (down), toward zero, and to nearest.
\end{definition}

\caution{The default rounding mode, which the list \enquote{round up,
round down} omits, is round to nearest with ties to even
\cite{ieee754-2019}, as in \cref{ex:l12-bf16}.}%
BF16, TF32 and FP8 are not IEEE 754 interchange formats; they borrow its
encoding rules, and E4M3 departs from them at the top of its
range.

\subsection{Quantization in practice}

Fewer bits mean fewer and cheaper operations, so ML typically
quantizes the inputs of an operation to a low precision and keeps a
high-precision accumulator.  The values of a tensor are rarely spread
evenly: most weights and gradients cluster near zero, with a long tail.
A format whose values are spaced evenly in the log domain then wastes
codes, and ML practice transforms the values before quantizing them, a
non-uniform quantization.\margin{The Transformer Engine of
\cref{sec:l12-te} applies the simplest such transform: it scales a tensor,
from statistics of its values, before casting it to FP8
\cite{nvidia2022hopper}.}  The resolution near zero matters most.

\begin{definition}[S2FP8]
  The format \term{shifted and squeezed FP8 (S2FP8)}[S2FP8]
  \cite{cambier2020s2fp8} stores a tensor $X$ in FP8 after a power
  transform with two per-tensor parameters $\alpha$ (squeeze) and $\beta$
  (shift):
  \begin{equation}
    X_{\text{truncated}} =
    \Bigl[\, 2^{-\beta} \cdot
      \mathrm{truncate}_{\text{FP8}}\bigl(2^{\beta} \cdot |X|^{\alpha}\bigr)
    \Bigr]^{1/\alpha} .
    \label{eq:l12-s2fp8}
  \end{equation}
  The transform acts on magnitudes; each element keeps the sign of $X$.
  In the log domain, $\log_2 (2^{\beta} |X|^{\alpha}) = \alpha \log_2 |X|
  + \beta$: $\alpha$ scales the spread of the exponents to fit the range
  of FP8 and $\beta$ shifts their centre, so the codes of the format fall
  where the values of the tensor are.  Both are computed for each tensor
  from the mean $\mu$ and the maximum $m$ of $\log_2 |X|$ over its
  non-zero elements: $\alpha = 15/(m - \mu)$ and $\beta = -\alpha\mu$, so
  that the transformed values have a mean exponent of~0 and a maximum
  of~15 \cite[§3.2]{cambier2020s2fp8}.
\end{definition}

\begin{keypoint}
  A floating-point format is a choice between range (exponent bits) and
  precision (fraction bits) at a given width.  ML lowers the width of
  the inputs and keeps the accumulator wide; each halving of the input
  width doubles the arithmetic intensity, and on recent tensor cores
  also the peak rate.
\end{keypoint}

\section{Tensor core}
\label{sec:l12-tensor-core}

A tensor core is one instance of a general task: adding an accelerator
for ML to a GPU.\source{M12, tensor core design; CPG §7.24,
§7.26, §7.27; PMPP ch.~14; \cite{nvidia2020ampere,nvidia2022hopper}.}
This section follows the design questions in order: which function to
build, how to build the compute units, how to store and move their
data, and how to exploit sparsity.

\subsection{Design questions and steps}

The design of an ML accelerator component raises four questions: what
functionality to add; what benefit it brings over the existing GPU;
how to build the compute units, and how many; and how to store and move
the data they need.  The answers follow a common sequence of steps.
\begin{enumerate}
  \item Identify the frequently executed operations, from knowledge of
    the algorithms or from profiling.
  \item Estimate the benefit over the existing data-parallel units:
    how much faster the operation would be in hardware than in
    software with the existing ISA.\margin{For the V100, at equal precision, a
    factor of about 4: \SI{125}{\tera\flop\per\second} on the FP16 tensor
    cores against \SI{31.4}{\tera\flop\per\second} for FP16 on the SIMT
    lanes (AWP Table~4 \cite{nvidia2020ampere}); against FP32's
    \SI{15.7}{\tera\flop\per\second} the factor of 8 includes the
    narrower format.}
  \item Design the interface: the new instructions, how data are fed to
    the unit, the storage format, and whether operands live in registers
    or in memory (private or shared).
  \item Decide whether to merge accelerator components into one unit,
    weighing programmability and usability against hardware cost and
    timing overhead.\margin{Step~3 is where the compiler of Lessons~10
    and~11 enters: a new unit is reached only through new instructions,
    which the compiler or a library must emit.}
\end{enumerate}

\subsection{The matrix multiply-accumulate unit}

GPUs chose an MMA unit because GEMM dominates ML
(\cref{sec:l12-gpu-ml}), and SIMD lanes still work one row--column pair
at a time.  The design space is between many small MMA units and one
large unit such as a systolic array, and every unit takes chip area
from other GPU functions; programmability matters as much as the raw
rate.  NVIDIA started with a $4 \times 4$ unit (Volta's
$4 \times 4 \times 4$ product per clock), and other vendors have
similar units: Intel's \term{Advanced Matrix Extensions (AMX)}[AMX]
add $16 \times 16$ matrix tiles and tile instructions to x86
CPUs.\caution{Lecture vs documentation: Intel documents
eight AMX tile registers of at most 16 rows of 64 bytes each; its INT8
example multiplies $16 \times 64$ inputs into a $16 \times 16$ INT32
result \cite{intel2024amxintrinsics}.  So $16 \times 16$ is the output
tile of that example, not a fixed matrix size.}

Within an MMA unit, three organisations are possible.
\begin{description}
  \item[Systolic array.]  Efficient for large tiles, but the fill and
    drain time is wasted on small ones: a $4 \times 4$ array was 40\,\%
    busy on a single product (\cref{tab:l12-systolic-trace}).
  \item[Option 1: parallel units.]  One dot-product unit
    (\cref{fig:l12-dot-unit}) per output element: 16 units of 4 FMAs
    compute a whole $4 \times 4 \times 4$ product at once
    (\cref{fig:l12-mma-options}, left).
  \item[Option 2: pipelining.]  Fewer units, shared over time: row~1
    times column~1, then row~1 times column~2, and so on, with the load,
    compute and store of successive pairs overlapped
    (\cref{fig:l12-mma-options}, right).
\end{description}
Parallel units give the highest throughput for the most area; a
pipelined unit trades throughput for area.\margin{Option~1 performs $16
\times 4 = 64$ FMAs per cycle, the rate of one Volta tensor core
(\cref{tab:l12-tensor-cores}).}

\begin{figure}[tb]
  \centering
  % Two ways to build a 4x4x4 matrix multiply-accumulate unit from the
% dot-product step of figure l12-dot-unit.  Option 1 (left): 16 dot-product
% units of 4 FMAs each, one per output element, all working in the same
% cycle (64 FMAs).  Option 2 (right): fewer units, pipelined; the row-column
% pairs (r1,c1), (r1,c2), (r1,c3), ... pass through load, compute and store
% stages one cycle apart.
% Redraws the outline figure m12l3_matrix.
% Usage: % Two ways to build a 4x4x4 matrix multiply-accumulate unit from the
% dot-product step of figure l12-dot-unit.  Option 1 (left): 16 dot-product
% units of 4 FMAs each, one per output element, all working in the same
% cycle (64 FMAs).  Option 2 (right): fewer units, pipelined; the row-column
% pairs (r1,c1), (r1,c2), (r1,c3), ... pass through load, compute and store
% stages one cycle apart.
% Redraws the outline figure m12l3_matrix.
% Usage: % Two ways to build a 4x4x4 matrix multiply-accumulate unit from the
% dot-product step of figure l12-dot-unit.  Option 1 (left): 16 dot-product
% units of 4 FMAs each, one per output element, all working in the same
% cycle (64 FMAs).  Option 2 (right): fewer units, pipelined; the row-column
% pairs (r1,c1), (r1,c2), (r1,c3), ... pass through load, compute and store
% stages one cycle apart.
% Redraws the outline figure m12l3_matrix.
% Usage: \input{figures/l12-mma-options}
% Width: body (120 mm); the drawing is about 114 mm wide.
\begin{tikzpicture}[font=\scriptsize\sffamily,
    du/.style={draw, fill=ln-accent!12, minimum width=8mm,
      minimum height=5.5mm, inner sep=0pt},
    st/.style={draw, minimum width=9mm, minimum height=4.6mm, inner sep=0pt,
      text height=1.6ex, text depth=0.4ex}]
  % option 1
  \node[anchor=west] at (0,8mm)
    {\iflangja{方式1：並列ユニット}{Option 1: parallel units}};
  \foreach \i in {1,...,4} {
    \foreach \j in {1,...,4} {
      \node[du] at ({(\j-1)*9.5mm+4mm},{-(\i-1)*7mm}) {$c_{\i\j}$};
    }
  }
  \node[anchor=north west, align=left, text=ln-muted] at (0,-24mm)
    {\iflangja{4 FMA × 16 要素 = 64 FMA/サイクル}{4 FMAs $\times$ 16 elements = 64 FMAs per cycle}};
  % option 2
  \begin{scope}[xshift=46mm]
    \node[anchor=west] at (0,8mm)
      {\iflangja{方式2：パイプライン化}{Option 2: pipelined unit}};
    \foreach \r/\lab in {0/{(r1,c1)},1/{(r1,c2)},2/{(r1,c3)},3/{(r1,c4)}} {
      \node[anchor=east] at (13mm,{-\r*6mm}) {\lab};
      \node[st, fill=ln-tip!20] at ({17.5mm+\r*9mm},{-\r*6mm})
        {\iflangja{ロード}{load}};
      \node[st, fill=ln-accent!18] at ({26.5mm+\r*9mm},{-\r*6mm})
        {\iflangja{演算}{comp}};
      \node[st, fill=ln-caution!22] at ({35.5mm+\r*9mm},{-\r*6mm})
        {\iflangja{ストア}{store}};
    }
    \draw[->, ln-muted] (13mm,-24mm) -- node[below] {\iflangja{サイクル}{cycle}} (68mm,-24mm);
  \end{scope}
\end{tikzpicture}

% Width: body (120 mm); the drawing is about 114 mm wide.
\begin{tikzpicture}[font=\scriptsize\sffamily,
    du/.style={draw, fill=ln-accent!12, minimum width=8mm,
      minimum height=5.5mm, inner sep=0pt},
    st/.style={draw, minimum width=9mm, minimum height=4.6mm, inner sep=0pt,
      text height=1.6ex, text depth=0.4ex}]
  % option 1
  \node[anchor=west] at (0,8mm)
    {\iflangja{方式1：並列ユニット}{Option 1: parallel units}};
  \foreach \i in {1,...,4} {
    \foreach \j in {1,...,4} {
      \node[du] at ({(\j-1)*9.5mm+4mm},{-(\i-1)*7mm}) {$c_{\i\j}$};
    }
  }
  \node[anchor=north west, align=left, text=ln-muted] at (0,-24mm)
    {\iflangja{4 FMA × 16 要素 = 64 FMA/サイクル}{4 FMAs $\times$ 16 elements = 64 FMAs per cycle}};
  % option 2
  \begin{scope}[xshift=46mm]
    \node[anchor=west] at (0,8mm)
      {\iflangja{方式2：パイプライン化}{Option 2: pipelined unit}};
    \foreach \r/\lab in {0/{(r1,c1)},1/{(r1,c2)},2/{(r1,c3)},3/{(r1,c4)}} {
      \node[anchor=east] at (13mm,{-\r*6mm}) {\lab};
      \node[st, fill=ln-tip!20] at ({17.5mm+\r*9mm},{-\r*6mm})
        {\iflangja{ロード}{load}};
      \node[st, fill=ln-accent!18] at ({26.5mm+\r*9mm},{-\r*6mm})
        {\iflangja{演算}{comp}};
      \node[st, fill=ln-caution!22] at ({35.5mm+\r*9mm},{-\r*6mm})
        {\iflangja{ストア}{store}};
    }
    \draw[->, ln-muted] (13mm,-24mm) -- node[below] {\iflangja{サイクル}{cycle}} (68mm,-24mm);
  \end{scope}
\end{tikzpicture}

% Width: body (120 mm); the drawing is about 114 mm wide.
\begin{tikzpicture}[font=\scriptsize\sffamily,
    du/.style={draw, fill=ln-accent!12, minimum width=8mm,
      minimum height=5.5mm, inner sep=0pt},
    st/.style={draw, minimum width=9mm, minimum height=4.6mm, inner sep=0pt,
      text height=1.6ex, text depth=0.4ex}]
  % option 1
  \node[anchor=west] at (0,8mm)
    {\iflangja{方式1：並列ユニット}{Option 1: parallel units}};
  \foreach \i in {1,...,4} {
    \foreach \j in {1,...,4} {
      \node[du] at ({(\j-1)*9.5mm+4mm},{-(\i-1)*7mm}) {$c_{\i\j}$};
    }
  }
  \node[anchor=north west, align=left, text=ln-muted] at (0,-24mm)
    {\iflangja{4 FMA × 16 要素 = 64 FMA/サイクル}{4 FMAs $\times$ 16 elements = 64 FMAs per cycle}};
  % option 2
  \begin{scope}[xshift=46mm]
    \node[anchor=west] at (0,8mm)
      {\iflangja{方式2：パイプライン化}{Option 2: pipelined unit}};
    \foreach \r/\lab in {0/{(r1,c1)},1/{(r1,c2)},2/{(r1,c3)},3/{(r1,c4)}} {
      \node[anchor=east] at (13mm,{-\r*6mm}) {\lab};
      \node[st, fill=ln-tip!20] at ({17.5mm+\r*9mm},{-\r*6mm})
        {\iflangja{ロード}{load}};
      \node[st, fill=ln-accent!18] at ({26.5mm+\r*9mm},{-\r*6mm})
        {\iflangja{演算}{comp}};
      \node[st, fill=ln-caution!22] at ({35.5mm+\r*9mm},{-\r*6mm})
        {\iflangja{ストア}{store}};
    }
    \draw[->, ln-muted] (13mm,-24mm) -- node[below] {\iflangja{サイクル}{cycle}} (68mm,-24mm);
  \end{scope}
\end{tikzpicture}

  \caption{Two ways to build a $4 \times 4 \times 4$ MMA unit.  Option 1:
    16 dot-product units, one per element of $C$.  Option 2: a pipelined
    unit through which the row--column pairs pass one cycle apart.}
  \label{fig:l12-mma-options}
\end{figure}

\subsection{Storing and moving the operands}

The operands of an MMA unit come from memory, and their volume is
large: a $4 \times 4 \times 4$ operation needs at least
$3 \times 16$ values ($A$, $B$ and $C$) in registers.  NVIDIA's
registers are private to a thread, which complicates feeding one unit
from them; the warp-level interface of \cref{sec:l12-gpu-ml} spreads
the operands over the registers of 32 threads.  The alternative is a
memory close to the unit, shared memory (Lesson~3), at the cost of
moving the data there from global memory first.\margin{Register file
size and banking, and the pressure that wide operands put on them, are
the subject of Lessons~4 and~6.}

For that move NVIDIA added an asynchronous copy inside the GPU.  The
A100's instruction \enquote{loads data directly from global memory into
shared memory, optionally bypassing L1 cache, and eliminating the need
for intermediate register file (RF) usage}
\cite{nvidia2020ampere}.\caution{Two different copies are called
asynchronous: \texttt{cudaMemcpyAsync} of Lesson~5 is a host-issued
transfer between host and device; the A100 copy is a device instruction
(\texttt{cp.async} in PTX, \texttt{memcpy\_async} in CUDA C++, CPG
§7.27) between global and shared memory.}  Because the threads no
longer wait for the data, they need a way to learn when it has
arrived.

\begin{definition}[Asynchronous barrier]
  An \term{asynchronous barrier}[非同期バリア] is a barrier split into
  an \emph{arrive} and a \emph{wait} operation: a thread signals its
  arrival, continues with independent work, and waits only when it
  needs the result (CPG §7.26 Asynchronous Barrier).  Combined with the
  asynchronous copy, a block can start the copy of the next tile, compute
  on the current one, and wait for the copy just before it is
  used.\margin{\texttt{\_\_syncthreads()} (Lesson~3) is the
  synchronous case: arrival and wait are one operation.  Hopper extends
  the barrier to count the bytes of TMA transfers (the asynchronous
  transaction barrier \cite{nvidia2022hopper}).}
\end{definition}

\subsection{Attention: data movement beyond GEMM}

The same concern, keeping fast units fed, drives the main kernel of
transformer models.  The operation called
\term{attention}[アテンション] computes, for
query, key and value matrices $Q$, $K$, $V$ of $N$ rows,
$\mathrm{softmax}(QK^{\mathsf T}/\sqrt{d})\,V$, where $d$ is the row
length; the intermediate $N \times N$ score matrix is large, and a
straightforward implementation writes it to HBM and reads it back.
\term{FlashAttention}[FlashAttention] is \enquote{an IO-aware exact
attention algorithm that uses tiling to reduce the number of memory
reads/writes between GPU high bandwidth memory (HBM) and GPU on-chip
SRAM} \cite{dao2022flashattention}: it computes the same result without
storing the full score matrix in HBM.\margin{FlashAttention-2
partitions the work better across thread blocks and warps,
\enquote{reaching 50-73\,\% of the theoretical maximum FLOPs/s on A100}
\cite{dao2023flashattention2}.}  Its argument is the one of
\cref{eq:l12-ai}: attention is limited by memory traffic, not by
arithmetic, so the gain comes from fewer bytes.

\subsection{Sparse operations}

Many matrices of HPC and graph workloads are mostly zeros, and ML
creates more of them on purpose.

\begin{definition}[Sparsity, sparse matrix, compressed format]
  The \term{sparsity}[スパース性] of a matrix or tensor is the fraction
  of its elements that are zero; a \term{sparse matrix}[疎行列] is one with enough
  zeros that storing and computing only the non-zeros pays.  A
  \term{compressed sparse format}[疎行列圧縮形式] stores only the
  non-zeros together with their positions.  The simplest is the
  coordinate format (COO), one (row, column, value) triple per
  non-zero.\margin{Compressed sparse row (CSR), the usual format for
  SpMV, replaces the row indices of COO by one offset per row (PMPP
  ch.~14).}
\end{definition}

\Cref{fig:l12-sparse}(a) and (b) show a $4 \times 4$ matrix with six
non-zeros in dense and COO form.  COO needs 18 numbers where the dense
form needs 16: the index overhead outweighs the saved zeros unless the
matrix is much sparser, and the irregular indices make hardware support
hard.  ML therefore restricts where the non-zeros may be.

\begin{figure}[tb]
  \centering
  % Three representations of the same 4x4 sparse matrix (six non-zeros, at
% most one per pair of columns): (a) dense storage, 16 values; (b) the
% coordinate (COO) format, one (row, column, value) triple per non-zero, 18
% numbers; (c) a toy structured-sparse format that keeps one value per pair
% plus one bit saying whether it came from the left (L) or right (R) entry
% of the pair, 8 values and 8 bits.
% Redraws the outline figures m12l3_sparse and m12l3_sparse2 with our own
% matrix.
% Usage: % Three representations of the same 4x4 sparse matrix (six non-zeros, at
% most one per pair of columns): (a) dense storage, 16 values; (b) the
% coordinate (COO) format, one (row, column, value) triple per non-zero, 18
% numbers; (c) a toy structured-sparse format that keeps one value per pair
% plus one bit saying whether it came from the left (L) or right (R) entry
% of the pair, 8 values and 8 bits.
% Redraws the outline figures m12l3_sparse and m12l3_sparse2 with our own
% matrix.
% Usage: % Three representations of the same 4x4 sparse matrix (six non-zeros, at
% most one per pair of columns): (a) dense storage, 16 values; (b) the
% coordinate (COO) format, one (row, column, value) triple per non-zero, 18
% numbers; (c) a toy structured-sparse format that keeps one value per pair
% plus one bit saying whether it came from the left (L) or right (R) entry
% of the pair, 8 values and 8 bits.
% Redraws the outline figures m12l3_sparse and m12l3_sparse2 with our own
% matrix.
% Usage: \input{figures/l12-sparse}
% Width: body (120 mm); the drawing is about 110 mm wide.
\begin{tikzpicture}[font=\scriptsize\sffamily,
    c/.style={draw=ln-muted, minimum size=5.4mm, inner sep=0pt},
    nz/.style={c, fill=ln-accent!20},
    >={Stealth[length=1.3mm]}]
  % (a) dense
  \node[anchor=west] at (0,6mm) {(a) \iflangja{密な格納}{dense}};
  \foreach \row [count=\i from 0] in {{0,4,0,2},{7,0,0,0},{5,0,0,1},{0,3,0,0}} {
    \foreach \v [count=\j from 0] in \row {
      \ifnum\v=0
        \node[c] at ({\j*5.4mm+2.7mm},{-\i*5.4mm}) {0};
      \else
        \node[nz] at ({\j*5.4mm+2.7mm},{-\i*5.4mm}) {\v};
      \fi
    }
  }
  \draw[ln-caution, thick] (10.8mm,3.2mm) -- (10.8mm,-19mm);
  % (b) COO
  \begin{scope}[xshift=34mm]
    \node[anchor=west] at (-2mm,6mm) {(b) COO};
    \node[text=ln-muted] at (2.5mm,1.5mm) {\iflangja{行}{row}};
    \node[text=ln-muted] at (9mm,1.5mm) {\iflangja{列}{col}};
    \node[text=ln-muted] at (15.5mm,1.5mm) {\iflangja{値}{val}};
    \foreach \r/\cc/\v [count=\n from 1] in {1/2/4,1/4/2,2/1/7,3/1/5,3/4/1,4/2/3} {
      \node at (2.5mm,{-\n*3.6mm+0.6mm}) {\r};
      \node at (9mm,{-\n*3.6mm+0.6mm}) {\cc};
      \node[text=ln-accent] at (15.5mm,{-\n*3.6mm+0.6mm}) {\v};
    }
  \end{scope}
  % (c) structured
  \begin{scope}[xshift=66mm]
    \node[anchor=west] at (0,6mm)
      {(c) \iflangja{1/2 構造化}{1-of-2 structured}};
    \foreach \row [count=\i from 0] in {{4/R,2/R},{7/L,0/L},{5/L,1/R},{3/R,0/L}} {
      \foreach \v/\b [count=\j from 0] in \row {
        \node[nz] at ({\j*13mm+2.7mm},{-\i*5.4mm}) {\v};
        \node[c, fill=ln-tip!22, minimum width=4mm] at ({\j*13mm+7.8mm},{-\i*5.4mm}) {\b};
      }
    }
    \node[anchor=north west, text=ln-muted, align=left] at (0,-19mm)
      {\iflangja{値8個 + 位置ビット8個}{8 values + 8 position bits}};
  \end{scope}
  \node[anchor=north west, text=ln-muted, align=left] at (0,-19mm)
    {\iflangja{値16個}{16 values}};
  \node[anchor=north west, text=ln-muted, align=left] at (32mm,-23.5mm)
    {\iflangja{18個の数}{18 numbers}};
\end{tikzpicture}

% Width: body (120 mm); the drawing is about 110 mm wide.
\begin{tikzpicture}[font=\scriptsize\sffamily,
    c/.style={draw=ln-muted, minimum size=5.4mm, inner sep=0pt},
    nz/.style={c, fill=ln-accent!20},
    >={Stealth[length=1.3mm]}]
  % (a) dense
  \node[anchor=west] at (0,6mm) {(a) \iflangja{密な格納}{dense}};
  \foreach \row [count=\i from 0] in {{0,4,0,2},{7,0,0,0},{5,0,0,1},{0,3,0,0}} {
    \foreach \v [count=\j from 0] in \row {
      \ifnum\v=0
        \node[c] at ({\j*5.4mm+2.7mm},{-\i*5.4mm}) {0};
      \else
        \node[nz] at ({\j*5.4mm+2.7mm},{-\i*5.4mm}) {\v};
      \fi
    }
  }
  \draw[ln-caution, thick] (10.8mm,3.2mm) -- (10.8mm,-19mm);
  % (b) COO
  \begin{scope}[xshift=34mm]
    \node[anchor=west] at (-2mm,6mm) {(b) COO};
    \node[text=ln-muted] at (2.5mm,1.5mm) {\iflangja{行}{row}};
    \node[text=ln-muted] at (9mm,1.5mm) {\iflangja{列}{col}};
    \node[text=ln-muted] at (15.5mm,1.5mm) {\iflangja{値}{val}};
    \foreach \r/\cc/\v [count=\n from 1] in {1/2/4,1/4/2,2/1/7,3/1/5,3/4/1,4/2/3} {
      \node at (2.5mm,{-\n*3.6mm+0.6mm}) {\r};
      \node at (9mm,{-\n*3.6mm+0.6mm}) {\cc};
      \node[text=ln-accent] at (15.5mm,{-\n*3.6mm+0.6mm}) {\v};
    }
  \end{scope}
  % (c) structured
  \begin{scope}[xshift=66mm]
    \node[anchor=west] at (0,6mm)
      {(c) \iflangja{1/2 構造化}{1-of-2 structured}};
    \foreach \row [count=\i from 0] in {{4/R,2/R},{7/L,0/L},{5/L,1/R},{3/R,0/L}} {
      \foreach \v/\b [count=\j from 0] in \row {
        \node[nz] at ({\j*13mm+2.7mm},{-\i*5.4mm}) {\v};
        \node[c, fill=ln-tip!22, minimum width=4mm] at ({\j*13mm+7.8mm},{-\i*5.4mm}) {\b};
      }
    }
    \node[anchor=north west, text=ln-muted, align=left] at (0,-19mm)
      {\iflangja{値8個 + 位置ビット8個}{8 values + 8 position bits}};
  \end{scope}
  \node[anchor=north west, text=ln-muted, align=left] at (0,-19mm)
    {\iflangja{値16個}{16 values}};
  \node[anchor=north west, text=ln-muted, align=left] at (32mm,-23.5mm)
    {\iflangja{18個の数}{18 numbers}};
\end{tikzpicture}

% Width: body (120 mm); the drawing is about 110 mm wide.
\begin{tikzpicture}[font=\scriptsize\sffamily,
    c/.style={draw=ln-muted, minimum size=5.4mm, inner sep=0pt},
    nz/.style={c, fill=ln-accent!20},
    >={Stealth[length=1.3mm]}]
  % (a) dense
  \node[anchor=west] at (0,6mm) {(a) \iflangja{密な格納}{dense}};
  \foreach \row [count=\i from 0] in {{0,4,0,2},{7,0,0,0},{5,0,0,1},{0,3,0,0}} {
    \foreach \v [count=\j from 0] in \row {
      \ifnum\v=0
        \node[c] at ({\j*5.4mm+2.7mm},{-\i*5.4mm}) {0};
      \else
        \node[nz] at ({\j*5.4mm+2.7mm},{-\i*5.4mm}) {\v};
      \fi
    }
  }
  \draw[ln-caution, thick] (10.8mm,3.2mm) -- (10.8mm,-19mm);
  % (b) COO
  \begin{scope}[xshift=34mm]
    \node[anchor=west] at (-2mm,6mm) {(b) COO};
    \node[text=ln-muted] at (2.5mm,1.5mm) {\iflangja{行}{row}};
    \node[text=ln-muted] at (9mm,1.5mm) {\iflangja{列}{col}};
    \node[text=ln-muted] at (15.5mm,1.5mm) {\iflangja{値}{val}};
    \foreach \r/\cc/\v [count=\n from 1] in {1/2/4,1/4/2,2/1/7,3/1/5,3/4/1,4/2/3} {
      \node at (2.5mm,{-\n*3.6mm+0.6mm}) {\r};
      \node at (9mm,{-\n*3.6mm+0.6mm}) {\cc};
      \node[text=ln-accent] at (15.5mm,{-\n*3.6mm+0.6mm}) {\v};
    }
  \end{scope}
  % (c) structured
  \begin{scope}[xshift=66mm]
    \node[anchor=west] at (0,6mm)
      {(c) \iflangja{1/2 構造化}{1-of-2 structured}};
    \foreach \row [count=\i from 0] in {{4/R,2/R},{7/L,0/L},{5/L,1/R},{3/R,0/L}} {
      \foreach \v/\b [count=\j from 0] in \row {
        \node[nz] at ({\j*13mm+2.7mm},{-\i*5.4mm}) {\v};
        \node[c, fill=ln-tip!22, minimum width=4mm] at ({\j*13mm+7.8mm},{-\i*5.4mm}) {\b};
      }
    }
    \node[anchor=north west, text=ln-muted, align=left] at (0,-19mm)
      {\iflangja{値8個 + 位置ビット8個}{8 values + 8 position bits}};
  \end{scope}
  \node[anchor=north west, text=ln-muted, align=left] at (0,-19mm)
    {\iflangja{値16個}{16 values}};
  \node[anchor=north west, text=ln-muted, align=left] at (32mm,-23.5mm)
    {\iflangja{18個の数}{18 numbers}};
\end{tikzpicture}

  \caption{One sparse matrix in three forms: (a) dense, 16 values; (b)
    COO triples, 18 numbers; (c) a structured format keeping one value
    per pair of columns (left of the red line, right of it) with a bit
    for its position, left (L) or right (R).}
  \label{fig:l12-sparse}
\end{figure}

\begin{definition}[Pruning, structured sparsity]
  The operation called \term{pruning}[枝刈り] sets weights of a trained
  network that contribute little to zero, which raises its sparsity.
  The constraint called \term{structured sparsity}[構造化スパース性]
  restricts the pattern:
  only a fixed fraction of the elements in every small group may be
  non-zero, so every group compresses to the same size and the index is
  short and regular.
\end{definition}

With 50\,\% structured sparsity and one non-zero per pair of elements,\caution{The one-bit left/right index here is a toy for one non-zero
in two.  In 2:4 sparsity each kept value can sit in any of four
positions, and for FP16 and BF16 the non-zeros of a group \enquote{are
indicated by two 2-bit indices in the metadata operand}
\cite{nvidia2026ptx} (\texttt{mma.sp}).}
16 elements are stored as 8 values plus 8 one-bit indices that say
whether each kept value was the left or the right element of its pair
(\cref{fig:l12-sparse}(c)).  The fixed layout simplifies the hardware
and nearly halves the storage, and the unit skips the multiplications by
zero, which raises its throughput.

\begin{definition}[2:4 structured sparsity]
  The scheme \term{2:4 structured sparsity}[2:4構造化スパース性], introduced with
  the A100, is \enquote{a new 2:4 sparse matrix definition that allows
  two non-zero values in every four-entry vector}
  \cite{nvidia2020ampere}.\margin{Row~1 of the matrix
of \cref{fig:l12-sparse}, $[0, 4, 0, 2]$, is already 2:4: it keeps
$4$ and $2$ with the indices $1$ and $3$ (counting from~0).}  A row of $A$ is stored as half its values
  plus metadata giving the positions of the two kept values in each
  group of four.
\end{definition}


The A100's sparse format \enquote{can be compressed efficiently and
reduce memory storage and bandwidth by almost 2x}
\cite{nvidia2020ampere},\tip{2:4 with FP16 stores a group of four
as two values plus two 2-bit indices: $32 + 4 = 36$ of 64 bits, about
56\,\%.} and its sparse tensor-core rates are twice the
dense ones (\cref{tab:l12-tensor-cores}).  Beyond
matrix--matrix products,
\term{sparse matrix-vector multiplication (SpMV)}[疎行列ベクトル積],
$y = Ax$ with a sparse $A$, is the core of many HPC and graph
algorithms, and accelerating it is an active research topic: its
intensity is low and its accesses irregular, so it profits little from
a dense matrix unit.

\begin{keypoint}[Lesson summary]
  GPUs suit ML: it is floating-point heavy, data parallel and tolerant
  of relaxed semantics, and GPUs have HBM and adopt new formats quickly.
  GEMM dominates DNNs; its intensity $MNK/(MK + NK + MN)$ (FP16) against
  the ops:byte ratio tells whether it is math- or memory-limited.  An
  MMA unit (the tensor core) replaces 16 SIMD dot products of a
  $4 \times 4$ product by one instruction; systolic arrays pay fill and
  drain.  A floating-point format trades exponent bits (range) for
  fraction bits (precision); ML narrows the inputs (BF16, TF32, FP8),
  keeps the accumulator wide and transforms values (S2FP8).  A
  tensor-core design chooses parallel or pipelined units, operand
  storage fed by asynchronous copies and barriers, and a sparse format:
  2:4 structured sparsity nearly halves storage and doubles the peak
  tensor-core rate.
\end{keypoint}

\end{document}
